\chapter{经典规范场与因果电磁响应}\label{chap:classical-gauge-response}

旋量表示给出了电子场怎样随时空变换，但尚未决定它在外场中怎样演化。作用量将变换规律与动力学联系起来：整体相位对称产生守恒电流，逐点改变相位基准则要求引入电磁规范势。

由同一作用量变分，电子电流成为 Maxwell 方程的源，规范势又进入 Dirac 方程。给定材料和边界以后，因果 Green 张量把局域方程变成源到场的响应。这一步先建立经典场的耦合与传播；场自身的量子涨落随后还需要扩大状态空间。

% between spacetime representation theory and local gauge symmetry.

\section{作用量、场方程与 Noether 守恒律}\label{sec:action-noether-dp61}

表示论告诉我们场可以怎样变换，却还没有规定哪一种运动由自然界实现。作用量把“允许的场表示”与“实际的动力学”连接起来，而 Noether 定理进一步把连续对称性转化为守恒量，为随后把整体相位对称局域化做好准备。

\subsection{从局域拉氏量到场方程}

质点力学里，Euler--Lagrange 方程来自“真实轨迹使作用量驻值”。场论没有改变这个原则，只是把有限个坐标 \(q_a(t)\) 换成每个时空点都有取值的场 \(\Phi_A(x)\)。因此这里首先要回答的不是“场方程长什么样”，而是怎样把普通变分法推广到局域场，并看清边界项为何决定运动方程与守恒流。

令 $\Phi_A(x)$ 表示一组经典场，$A$ 只用于枚举不同场分量。全文取 $x^0=ct$，因此 $\dd^4 x=c\,\dd t\,\dd^3 r$。若拉氏量密度只依赖同一点的场及其一阶导数，作用量定义为
\begin{equation}
 {
 S[\Phi]=\frac1c\int_{\Omega}\dd^4 x\,
 \Lag\!\left(\Phi_A,\partial_\mu\Phi_A\right).}
 \label{eq:field-action-general-dp61}
\end{equation}
这里 $\Omega$ 是所考察的时空区域；$[\Lag]=\mathrm{J\,m^{-3}}$，所以 $[S]=\mathrm{J\,s}$。作用量 $S$ 是把整段时空历史映射到一个具有作用量量纲标量的泛函；其驻值条件决定经典运动方程。

对场作任意小变分 $\Phi_A\to\Phi_A+\delta\Phi_A$。只要在边界 $\partial\Omega$ 上固定场，即 $\delta\Phi_A|_{\partial\Omega}=0$，驻值条件 $\delta S=0$ 就要求
\begin{equation}
 {
 \partial_\mu\frac{\partial\Lag}{\partial(\partial_\mu\Phi_A)}
 -\frac{\partial\Lag}{\partial\Phi_A}=0.}
 \label{eq:field-euler-lagrange-dp61}
\end{equation}
这是场的 Euler--Lagrange 方程。它与质点力学中的 Euler--Lagrange 方程完全同源；差别只在于离散自由度指标被时空点和场分量共同替代。

\subsection{连续对称怎样变成守恒流}

场方程告诉我们给定初态以后怎样演化，但还没有解释为什么某些组合在整个演化中保持不变。连续对称性的作用正是把“变换后物理不变”转化成一个局域连续性方程。与只记住“对称性对应守恒量”相比，Noether 定理更强：它直接给出守恒流的表达式，并说明守恒荷只能通过边界流动而不能在体内凭空产生。

先考虑不移动时空坐标的连续内部变换。令 $\epsilon$ 是一个与 $x$ 无关的无穷小常数，并写
\begin{equation*}
 \delta\Phi_A=\epsilon\,\Delta_A(\Phi).
\end{equation*}
若拉氏量的变化最多是一个四维散度，
\begin{equation}
 \delta\Lag=\epsilon\,\partial_\mu K^\mu,
 \label{eq:noether-total-derivative-dp61}
\end{equation}
则作用量在固定边界条件下不变。利用 Euler--Lagrange 方程把体项消去以后，剩余变化必须写成局域连续性方程
\begin{equation}
 {
 \partial_\mu j^\mu=0,
 \qquad
 j^\mu
 =\sum_A
 \frac{\partial\Lag}{\partial(\partial_\mu\Phi_A)}\,\Delta_A
 -K^\mu.}
 \label{eq:noether-current-general-dp61}
\end{equation}
这就是固定坐标内部连续对称性的 Noether 定理。\eqnrefs{eq:noether-current-general-dp61} 比“对称性对应守恒量”更具体：它告诉我们守恒流怎样从拉氏量逐项计算出来。

若把该流写成 $j^\mu=(c\rho_j,\vb{j})$，则
\begin{equation*}
 \partial_\mu j^\mu=0
 \quad\Longleftrightarrow\quad
 \partial_t\rho_j+\vb{\nabla}\!\cdot\vb{j}=0.
\end{equation*}
对任意固定空间体积 $V$ 积分并用 Gauss 定理，得到
\begin{equation}
 {
 \frac{\dd}{\dd t}\int_V\dd^3 r\,\rho_j
 =-\oint_{\partial V}\vb{j}\cdot\dd\vb{S}.}
 \label{eq:noether-flux-balance-dp61}
\end{equation}
左边是区域内守恒荷的变化率，右边是穿过边界的净流出。若取全空间且场在无穷远处衰减得足够快，边界通量消失，于是总荷
\begin{equation*}
 Q_j=\int\dd^3 r\,\rho_j
\end{equation*}
不随时间变化。局域连续性方程因此比“总量守恒”更强：它还规定守恒量只能通过边界输运，不能在体内凭空产生或消失。

\begin{exbox}[自由 Dirac 场的整体 $U(1)$]
自由 Dirac 拉氏量为
\begin{equation*}
 \Lag_D
 =\bar\psi\left(\ii\hbar c\gamma^\mu\partial_\mu-m_{\mathrm e} c^2\right)\psi.
\end{equation*}
它在常数相位变换 $\psi\to \ee^{\ii\theta}\psi$、$\bar\psi\to\bar\psi \ee^{-\ii\theta}$ 下不变。求粒子数流、电流及其守恒荷，核对时间分量与概率密度的关系。
\textbf{解。} Noether 流的整体倍数取决于生成元的归一化；选择粒子数流的常用归一化后，
\begin{equation}
 {n^\mu=c\,\bar\psi\gamma^\mu\psi,
 \qquad \partial_\mu n^\mu=0.}
 \label{eq:dirac-number-current-dp61}
\end{equation}
于是
\begin{equation*}
 n^0=c\psi^\dagger\psi,
 \qquad
 N=\int\dd^3 r\,\psi^\dagger\psi.
\end{equation*}
若该 Dirac 场携带电荷 $q_{\mathrm e}$，物理电流就是
\begin{equation}
 {j^\mu=q_{\mathrm e} n^\mu=q_{\mathrm e} c\,\bar\psi\gamma^\mu\psi.}
 \label{eq:dirac-electric-current-from-noether-dp61}
\end{equation}
整体 $U(1)$ 对称取 $\theta$ 为时空常数，因此普通导数不产生额外项。令相位依赖时空位置 $\theta=\theta(x)$ 后，$\partial_\mu\psi$ 会额外产生与 $\partial_\mu\theta$ 成正比的项。抵消该项需要引入电磁四势，并用协变导数取代普通导数。对应关系为
\[
 \text{整体相位对称性}
 \longrightarrow \text{Noether 守恒流}
 \longrightarrow \text{局域化相位}
 \longrightarrow \text{规范联络 }A_\mu.
\]

\medskip

一般 Noether 流~\eqnrefs{eq:noether-current-general-dp61} 已经建立。对自由 Dirac 拉氏量，只需把整体相位变换 $\delta\psi=\ii\epsilon\psi$、$\delta\bar\psi=-\ii\epsilon\bar\psi$ 代入该公式。质量项与动能项的相位变化彼此抵消，所以 $K^\mu=0$；再按“每个粒子携带单位数荷”的约定重标度生成元，就得到\eqnrefs{eq:dirac-number-current-dp61}。

取无穷小整体变换
\begin{equation*}
 \delta\psi=\ii\epsilon\psi,
 \qquad
 \delta\bar\psi=-\ii\epsilon\bar\psi,
 \qquad \partial_\mu\epsilon=0.
\end{equation*}
质量项的变化为
\begin{equation*}
 \delta(\bar\psi\psi)
 =(-\ii\epsilon\bar\psi)\psi
 +\bar\psi(\ii\epsilon\psi)=0.
\end{equation*}
动能项中
\begin{align*}
 \delta\!\left(\bar\psi\gamma^\mu\partial_\mu\psi\right)
 &=(-\ii\epsilon\bar\psi)\gamma^\mu\partial_\mu\psi
 +\bar\psi\gamma^\mu\partial_\mu(\ii\epsilon\psi),\\
 &=-\ii\epsilon\bar\psi\gamma^\mu\partial_\mu\psi
 +\ii\epsilon\bar\psi\gamma^\mu\partial_\mu\psi=0,
\end{align*}
所以 $K^\mu=0$。又因为
\begin{equation*}
 \frac{\partial\Lag_D}{\partial(\partial_\mu\psi)}
 =\ii\hbar c\,\bar\psi\gamma^\mu,
 \qquad
 \frac{\partial\Lag_D}{\partial(\partial_\mu\bar\psi)}=0,
\end{equation*}
Noether 公式给出的原始流为
\begin{equation*}
 j_{\rm raw}^\mu
 =(\ii\hbar c\bar\psi\gamma^\mu)(\ii\psi)
 =-\hbar c\bar\psi\gamma^\mu\psi.
\end{equation*}
连续群生成元可以整体重标度；相应 Noether 流也按同一常数重标度。把生成元归一化成“每个粒子携带单位数荷”，定义
\begin{equation*}
 n^\mu\equiv-\frac1\hbar j_{\rm raw}^\mu
 =c\bar\psi\gamma^\mu\psi.
\end{equation*}
于是 $\partial_\mu n^\mu=0$。时间分量使用 $\bar\psi=\psi^\dagger\gamma^0$：
\begin{equation*}
 n^0=c\bar\psi\gamma^0\psi
 =c\psi^\dagger\psi.
\end{equation*}
故数荷为 $N=c^{-1}\int\dd^3 r\,n^0=\int\dd^3 r\,\psi^\dagger\psi$。
\end{exbox}


\begin{theorem}[一阶局部作用量的 Euler--Lagrange 方程]
对 $S[\phi]=\int\mathcal L(\phi,\partial_\mu\phi,x)\,\dd^dx$，若变分在边界消失，则驻值条件等价于 $\partial\mathcal L/\partial\phi_a-\partial_\mu(\partial\mathcal L/\partial(\partial_\mu\phi_a))=0$。复场的共轭分量按独立实自由度处理。
\end{theorem}

\begin{proof}
正文\eqnrefs{eq:field-action-general-dp61} 把场的整段时空历史映射成作用量。对场做任意小变分时，唯一需要处理的是含 $\partial_\mu\delta\Phi_A$ 的项；把它分部积分后，边界条件 $\delta\Phi_A|_{\partial\Omega}=0$ 消去表面项，而区域内部的 $\delta\Phi_A(x)$ 又可以独立任取，因此体积分系数必须逐点为零。这正给出正文\eqnrefs{eq:field-euler-lagrange-dp61}。

从定义出发，
\begin{align*}
 \delta S
 &=\frac1c\int_\Omega\dd^4 x\,
 \left[
 \frac{\partial\Lag}{\partial\Phi_A}\delta\Phi_A
 +\frac{\partial\Lag}{\partial(\partial_\mu\Phi_A)}
 \delta(\partial_\mu\Phi_A)
 \right],\\
 &=\frac1c\int_\Omega\dd^4 x\,
 \left[
 \frac{\partial\Lag}{\partial\Phi_A}\delta\Phi_A
 +\frac{\partial\Lag}{\partial(\partial_\mu\Phi_A)}
 \partial_\mu\delta\Phi_A
 \right].
\end{align*}
第二项使用乘积求导反解：
\begin{align*}
 \frac{\partial\Lag}{\partial(\partial_\mu\Phi_A)}
 \partial_\mu\delta\Phi_A
 &=\partial_\mu\!\left[
 \frac{\partial\Lag}{\partial(\partial_\mu\Phi_A)}
 \delta\Phi_A\right]
 -\partial_\mu\!\left[
 \frac{\partial\Lag}{\partial(\partial_\mu\Phi_A)}\right]
 \delta\Phi_A.
\end{align*}
代回以后，
\begin{align*}
 \delta S
 &=\frac1c\int_\Omega\dd^4 x\,
 \left[
 \frac{\partial\Lag}{\partial\Phi_A}
 -\partial_\mu\frac{\partial\Lag}{\partial(\partial_\mu\Phi_A)}
 \right]\delta\Phi_A\\
 &\quad+\frac1c\int_\Omega\dd^4 x\,
 \partial_\mu\!\left[
 \frac{\partial\Lag}{\partial(\partial_\mu\Phi_A)}
 \delta\Phi_A\right].
\end{align*}
四维 Gauss 定理把最后一项变成 $\partial\Omega$ 上的面积分；由于边界上 $\delta\Phi_A=0$，该项消失。于是
\begin{equation*}
 \delta S
 =\frac1c\int_\Omega\dd^4 x\,
 \mathfrak E_A(x)\,\delta\Phi_A(x),
\end{equation*}
其中
\begin{equation*}
 \mathfrak E_A
 \equiv
 \frac{\partial\Lag}{\partial\Phi_A}
 -\partial_\mu\frac{\partial\Lag}{\partial(\partial_\mu\Phi_A)}.
\end{equation*}
区域内部的 $\delta\Phi_A(x)$ 可以独立任取；要使 $\delta S=0$ 对所有变分都成立，只能有 $\mathfrak E_A(x)=0$。移项即得到\eqnrefs{eq:field-euler-lagrange-dp61}。
\end{proof}

\begin{theorem}[Noether 流的在壳守恒]
若无穷小场变换使拉氏量密度变化为全导数 $\delta\mathcal L=\partial_\mu K^\mu$，则 $j^\mu=\frac{\partial\mathcal L}{\partial(\partial_\mu\phi_a)}\delta\phi_a-K^\mu$ 在 Euler--Lagrange 方程成立时满足 $\partial_\mu j^\mu=0$；坐标变换需使用本章的固定坐标总变分。
\end{theorem}

\begin{proof}
链式法则给

\begin{align*}
 \delta\Lag
 &=\frac{\partial\Lag}{\partial\Phi_A}\delta\Phi_A
 +\frac{\partial\Lag}{\partial(\partial_\mu\Phi_A)}
 \partial_\mu\delta\Phi_A.
\end{align*}
对第二项做同样的乘积拆分，
\begin{align*}
 \delta\Lag
 &=\left[
 \frac{\partial\Lag}{\partial\Phi_A}
 -\partial_\mu\frac{\partial\Lag}{\partial(\partial_\mu\Phi_A)}
 \right]\delta\Phi_A\\
 &\quad+\partial_\mu\!\left[
 \frac{\partial\Lag}{\partial(\partial_\mu\Phi_A)}
 \delta\Phi_A\right].
\end{align*}
在满足 Euler--Lagrange 方程的经典解上，第一行的方括号为零。再代入 $\delta\Phi_A=\epsilon\Delta_A$：
\begin{equation*}
 \delta\Lag
 =\epsilon\,\partial_\mu\!\left[
 \sum_A\frac{\partial\Lag}{\partial(\partial_\mu\Phi_A)}
 \Delta_A\right].
\end{equation*}
另一方面，对称性假设给 $\delta\Lag=\epsilon\partial_\mu K^\mu$。两式相减，并利用常数 $\epsilon$ 可约去，得到
\begin{equation*}
 \partial_\mu\!\left[
 \sum_A\frac{\partial\Lag}{\partial(\partial_\mu\Phi_A)}
 \Delta_A-K^\mu\right]=0,
\end{equation*}
即\eqnrefs{eq:noether-current-general-dp61}。
\end{proof}


\section{局域 $U(1)$ 与 Maxwell--Dirac 理论}

时空表示已经确定场怎样随坐标变换；带电场还具有独立的内部相位自由度。Noether 分析说明整体 $U(1)$ 对称对应守恒电荷，但若相位基准允许随时空位置变化，普通导数会额外作用在局域相位上而失去协变性。补偿这一项所需的规范联络就是电磁四势。把协变导数写入作用量以后，Maxwell 方程与带电 Dirac 方程从同一个局域 $U(1)$ 原则得到；哈密顿量形式随后自然区分正则动量与动力学动量。

\subsection{协变导数}

自由 Dirac 场具有全局 $U(1)$ 相位对称性。若希望相位选择可以逐点改变，普通导数就会感受到局域相位的梯度，因此必须引入一个能够补偿这项变化的联络。采用全书约定，带电场在局域变换下满足
\begin{equation*}
\psi'(x)=\exp\!\left(\frac{i q_{\mathrm e}}{\hbar}\chi(x)\right)\psi(x),
\qquad q_{\mathrm e}=-e,
\end{equation*}
其中 $\chi$ 是任意实函数。要求导数项与 $\psi$ 具有相同的局域相位变换律，就必须定义
\begin{equation}
\boxed{D_\mu\equiv\partial_\mu+\frac{i q_{\mathrm e}}{\hbar}A_\mu,}
\label{eq:covariant-derivative}
\end{equation}
并要求规范势变换为
\begin{equation}
A'_\mu=A_\mu-\partial_\mu\chi.
\label{eq:gauge-A-lower}
\end{equation}
于是 $D'_\mu\psi'=\ee^{\ii q_{\mathrm e}\chi/\hbar}D_\mu\psi$。这就是“协变”二字的具体含义：$D_\mu\psi$ 与 $\psi$ 属于同一个局域表示，因此可以把自由 Dirac 项中的 $\partial_\mu$ 替换为 $D_\mu$ 而不破坏局域 $U(1)$ 对称性。

在三维记号中
\begin{equation*}
A^\mu=(\phi/c,\mathbf A),
\qquad
A_\mu=(\phi/c,-\mathbf A),
\end{equation*}
\eqnrefs{eq:gauge-A-lower} 等价于 $\mathbf A' =\mathbf A+\nabla\chi$、$\phi'=\phi-\partial_t\chi$。因此标势和矢势不是两个独立的“物理场”，而是同一个四维联络的分量。

联络的局域曲率由协变导数的对易子测量：
\begin{equation}
\boxed{
[D_\mu,D_\nu]=\frac{i q_{\mathrm e}}{\hbar}F_{\mu\nu},
\qquad
F_{\mu\nu}=\partial_\mu A_\nu-\partial_\nu A_\mu.}
\label{eq:covariant-commutator}
\end{equation}
$F_{\mu\nu}$ 在\eqnrefs{eq:gauge-A-lower} 下不变，其时空分量分别给电场与磁场。局域 $U(1)$ 因而把“为什么需要电磁势”和“为什么物理场强是势的反对称导数”连接成同一个几何结构。

\subsection{Maxwell--Dirac 方程}

局域 $U(1)$ 已经确定带电 Dirac 场必须通过 $D_\mu=\partial_\mu+\ii q_{\mathrm e} A_\mu/\hbar$ 与四势耦合。由同一作用量变分，可以同时得到物质方程、规范场方程和守恒流；再给波动算符指定因果边界条件，就得到 Green 张量与 QED 传播子共同使用的逆算符结构。

在平直时空中，由 $F_{\mu\nu}=\partial_\mu A_\nu-\partial_\nu A_\mu$ 构造的一阶导数、Lorentz 与局域 $U(1)$ 不变的最低阶规范动能是 $F_{\mu\nu}F^{\mu\nu}$。另一个二次标量 $F_{\mu\nu}{}^\star\!F^{\mu\nu}$ 是全导数，因此不改变这里讨论的局域 Euler--Lagrange 方程。电子与电磁场的最小耦合作用量于是写成
\begin{equation}
 \boxed{
 S[\psi,\bar\psi,A]
 =\frac1c\int\dd^4 x\,
 \left[
 \bar\psi\left(\ii\hbar c\gamma^\mu D_\mu-m_{\mathrm e} c^2\right)\psi
 -\frac{1}{4\mu_0}F_{\mu\nu}F^{\mu\nu}
 \right].}
 \label{eq:maxwell-dirac-action}
\end{equation}
展开协变导数，物质与规范场之间唯一的最小三点耦合为
\begin{equation}
 \boxed{
 \Lag_{\rm int}
 =-q_{\mathrm e} c\,\bar\psi\gamma^\mu\psi A_\mu.}
 \label{eq:qed-interaction-lagrangian}
\end{equation}
这个顶角随后既生成 QED 的微扰图，也生成自由电子--经典光场和自由电子--量子光场的低能约化。

把第一章已经建立的场论 Euler--Lagrange 方程分别用于 $\bar\psi$ 和 $A_\nu$，得到
\begin{equation}
 \boxed{
 \left(\ii\hbar c\slashed D-m_{\mathrm e} c^2\right)\psi=0,}
 \label{eq:dirac-em}
\end{equation}
以及
\begin{equation}
 \boxed{
 \partial_\mu F^{\mu\nu}=\mu_0 j^\nu,
 \qquad
 j^\nu\equiv q_{\mathrm e} c\,\bar\psi\gamma^\nu\psi.}
 \label{eq:maxwell-source}
\end{equation}
这里 $j^\mu$ 是物质场的电磁 Noether 流。全局相位对称性，等价地 Dirac 方程和伴随方程，给出
\begin{equation}
 \boxed{\partial_\mu j^\mu=0.}
 \label{eq:current-conservation}
\end{equation}
因此 Maxwell 源方程与物质动力学并不是两套独立输入：反对称性使 $\partial_\nu\partial_\mu F^{\mu\nu}=0$，而物质方程恰好保证相同的连续性条件。

后续介质和开放几何真正需要的是“给定源以后怎样选择因果解”。在 Lorenz 规范 $\partial_\mu A^\mu=0$ 下，熟悉的 Maxwell 方程压缩为
\begin{equation*}
 \Box A^\mu=\mu_0j^\mu,
 \qquad
 \Box\equiv\partial_\mu\partial^\mu.
\end{equation*}
给波动算符指定推迟边界条件，就定义其推迟逆核 $G_R^{(0)}$；真空中该逆核把源卷积成标准推迟势。对电子--场问题真正需要继续保留的是同一个波动算符、源项与因果逆核之间的关系：
\begin{equation*}
 A^\mu(x)=\mu_0\int\dd^4 x'\,G_R^{(0)}(x-x')j^\mu(x'),
 \qquad
 G_R^{(0)}(x-x')=0\quad(t<t').
\end{equation*}
Maxwell Green 张量把同一“带因果边界条件的逆算符”推广到色散介质、边界与开放几何；QED 传播子则把这一逆核思想提升到量子相关函数。

\subsection{正则动量与动力学动量}\label{sec:feqo-minimal-coupling}

从 \eqnrefs{eq:dirac-em} 出发，把时间导数单独提出。定义
\begin{equation*}
 \beta_D\equiv\gamma^0,
 \qquad
 \alpha^i\equiv\gamma^0\gamma^i,
 \qquad
 \hat{\mathbf p}\equiv-\ii\hbar\nabla.
\end{equation*}
逐项展开
\begin{align*}
 \ii\hbar c\gamma^\mu D_\mu
 &=\ii\hbar\gamma^0\partial_t
 +\ii\hbar c\gamma^i\partial_i
 -q_{\mathrm e}\gamma^0\phi
 +q_{\mathrm e} c\gamma^i A_i,
\end{align*}
将空间项与质量项移到右侧并左乘 $\gamma^0$，得到
\begin{equation}
 \boxed{
 \ii\hbar\partial_t\psi
 =\hat H_D\psi,
 \qquad
 \hat H_D
 =c\vb{\alpha}\cdot\hat{\vb{\Pi}}
 +\beta_D m_{\mathrm e} c^2+q_{\mathrm e}\phi,
 }
 \label{eq:field-dirac-hamiltonian-em}\end{equation}
其中
\begin{equation}
 \hat{\vb{\Pi}}\equiv\hat{\mathbf p}-q_{\mathrm e}\mathbf A
 \label{eq:kinetic-momentum}
\end{equation}
是动力学动量。$\hat{\mathbf p}$ 是与坐标正则共轭的正则动量；在有矢势时二者不能混同。

若外场依赖时间，$\hat H_D(t)$ 在不同时刻一般不对易。定义演化算符
\begin{equation*}
 \ii\hbar\partial_t\hat U_D(t,t_0)=\hat H_D(t)\hat U_D(t,t_0),
 \qquad \hat U_D(t_0,t_0)=\id().
\end{equation*}
逐次积分得到
\begin{equation}
 \hat U_D(t,t_0)=
 \Tord\exp\!\left[-\frac{\ii}{\hbar}
 \int_{t_0}^{t}\dd t'\,\hat H_D(t')\right].
 \label{eq:dirac-time-evolution}
\end{equation}
因此在进入自由极限之前，单电子外场问题已经由 $\hat H_D[A_\mu]$ 与 \eqnrefs{eq:dirac-time-evolution} 完整定义。相互作用绘景、程函相位和 PINEM 相位调制都可看成这个演化问题在不同表示、投影和尺度极限下的约化。

\begin{proposition}[跃迁电流的连续性与规范端点项]
设 $\Psi_i,\Psi_f$ 都满足同一实外势下的 Dirac 方程。定义复数值跃迁源
\begin{equation*}
 \rho_{fi}^{\rm ch}=q_{\mathrm e}\Psi_f^\dagger\Psi_i,
 \qquad \vb{J}_{fi}^{\rm ch}=q_{\mathrm e}c\Psi_f^\dagger\boldsymbol\alpha\Psi_i.
\end{equation*}
它们满足
\begin{equation}
 \partial_t\rho_{fi}^{\rm ch}+\nabla\cdot\vb{J}_{fi}^{\rm ch}=0.
 \label{eq:feqo-transition-current-continuity}
\end{equation}
若空间边界项消失，则线性耦合矩阵元
$V_{fi}=\int(\rho_{fi}^{\rm ch}\phi-\vb{J}_{fi}^{\rm ch}\cdot\vb{A})\dd^3r$
在 $\vb{A}' =\vb{A}+\nabla\chi$、$\phi'=\phi-\partial_t\chi$ 下满足
\begin{equation}
 \delta V_{fi}(t)=-\frac{\dd}{\dd t}\int\dd^3r\,\rho_{fi}^{\rm ch}(\vb{r},t)\chi(\vb{r},t).
 \label{eq:feqo-transition-gauge-endpoint}
\end{equation}
完整跃迁概率保持规范不变；有限时间散射必须连同初末态的规范相位处理。
\end{proposition}

\begin{proof}
分别将 $\Psi_i$ 的哈密顿量方程左乘 $\Psi_f^\dagger$，以及 $\Psi_f$ 方程的厄米共轭右乘 $\Psi_i$。相加得到
\begin{align*}
 \partial_t(\Psi_f^\dagger\Psi_i)
 &=-c\big[(\partial_j\Psi_f^\dagger)\alpha_j\Psi_i
 +\Psi_f^\dagger\alpha_j\partial_j\Psi_i\big]\\
 &=-\nabla\cdot(c\Psi_f^\dagger\boldsymbol\alpha\Psi_i),
\end{align*}
因为质量矩阵和同一外势的两个局域乘法项逐项抵消。再对规范变化的矩阵元作空间分部积分：
\begin{align*}
 \delta V_{fi}
 &=-\int\dd^3r\,[\rho_{fi}^{\rm ch}\partial_t\chi
 +\vb{J}_{fi}^{\rm ch}\cdot\nabla\chi]\\
 &=-\int\dd^3r\,[\rho_{fi}^{\rm ch}\partial_t\chi
 +(\partial_t\rho_{fi}^{\rm ch})\chi].
\end{align*}
这正是总时间导数。自由定态的时间因子为 $\ee^{-\ii\omega_{if}t}$，$\omega_{if}=(E_i-E_f)/\hbar$，故频率分量上 $\nabla\cdot\vb{J}_{fi}^{\rm ch}=\ii\omega_{if}\rho_{fi}^{\rm ch}$。纵向跃迁电流与电荷密度必须满足这一约束。
\end{proof}

\begin{exbox}[正负能隙与有限脉冲的混合振幅]
只考察一个已指定的正负能通道，令自由能量差 $\Delta=E_f-E_i$，$|\Delta|\ge2m_{\mathrm e}c^2$，外势矩阵元为 $v_{fi}(t)$。在有限脉冲和一阶微扰条件下，求正负能通道混合振幅的界，并判断大能隙是否足以压低混合。
\textbf{解。} 在 $\eta=\hbar^{-1}\int\|V(t)\|\dd t\ll1$ 的有限时间微扰区间，一阶混合振幅为
\begin{equation*}
 c_{fi}^{(1)}=-\frac{\ii}{\hbar}\int_{t_i}^{t_f}\dd t\,
 v_{fi}(t)\ee^{\ii\Delta t/\hbar},
 \qquad c_{fi}=c_{fi}^{(1)}+O(\eta^2).
\end{equation*}
假定 $v_{fi}$ 绝对连续，分部积分给
\begin{equation*}
 |c_{fi}^{(1)}|\le
 \frac{|v_{fi}(t_i)|+|v_{fi}(t_f)|
 +\int_{t_i}^{t_f}|\dot v_{fi}(t)|\dd t}{|\Delta|}.
\end{equation*}
对于端点关闭、$v_{fi}(t)=v_*w(t/T)$ 的光滑脉冲，分子变为 $|v_*|\int|w'(s)|\dd s$。如果谱在能隙附近共振，$\dot v_{fi}$ 则包含该快速振荡，单纯比较峰值 $|v_*|$ 与能隙不再够用。低光子能量、外场在 Compton 尺度上缓变、有限脉冲混合小，是不同的控制条件；强场多光子与隧穿对产生不能仅由 $\hbar\omega<2m_{\mathrm e}c^2$ 排除。
\end{exbox}


动力学动量包含位置依赖的矢势，因此不同方向的动力学平移一般不再对易。它们的对易子由磁场唯一确定：
\begin{equation}
 \boxed{
 [\hat\Pi_i,\hat\Pi_j]
 =\ii q_{\mathrm e}\hbar\epsilon_{ijk}B_k.}
 \label{eq:pi-commutator}
\end{equation}
这里 $i,j,k\in\{1,2,3\}$，$\epsilon_{ijk}$ 是三维 Levi--Civita 符号。\eqnrefs{eq:pi-commutator} 是协变导数曲率关系在空间分量上的单粒子形式：磁场同时控制经典 Lorentz 力与量子动力学平移的非对易性。

同一个曲率还出现在 Dirac 算符的平方中。将一阶方程 \eqnrefs{eq:dirac-em} 与共轭的一阶因子组合后，二阶方程为
\begin{equation}
 \boxed{
 \left[
 \hbar^2D_\mu D^\mu+m_{\mathrm e}^2c^2
 +\frac{q_{\mathrm e}\hbar}{2}\sigma^{\mu\nu}F_{\mu\nu}
 \right]\psi=0.
 }
 \label{eq:dirac-squared}
\end{equation}
前三个对象的来源彼此独立而在这一式中汇合：$D_\mu D^\mu$ 是规范协变的二阶平移，$m_{\mathrm e}^2c^2$ 来自相对论质量壳，$\sigma^{\mu\nu}F_{\mu\nu}$ 则把旋量 Lorentz 生成元与电磁曲率收缩。其低能展开将产生 Pauli 磁矩、Darwin 项和自旋--轨道耦合。

四分量 Dirac 理论的低能正能约化首先写成\cref{chap:math-foundations}建立的 $P/Q$ 消元问题；级数展开只在谱隙和算符范数给出控制参数之后进行。定义
\begin{equation}
 P_+\equiv\frac{\id()+\beta_D}{2},
 \qquad
 Q_-\equiv\frac{\id()-\beta_D}{2},
 \qquad
 P_++Q_-=\id(),
 \label{eq:dirac-positive-negative-projectors-r16}
\end{equation}
其中 $P_+$ 与 $Q_-$ 分别投影到 $\beta_D=+1$ 与 $\beta_D=-1$ 的静止能量二分量块。除 $\mathbf p=0$ 外，它们并不是自由 Dirac 哈密顿量的精确正、负能谱投影，因为 $c\boldsymbol\alpha\cdot\mathbf p$ 会混合两个块；低能正能电子支是在消去 $Q_-$ 块之后得到的有效谱支。对静态外场写
\begin{equation}
 \hat H_D=\beta_Dm_{\mathrm e} c^2+\hat{\mathcal V}+\hat{\mathcal O},
 \qquad
 \hat{\mathcal V}\equiv q_{\mathrm e}\phi,
 \qquad
 \hat{\mathcal O}\equiv c\boldsymbol\alpha\cdot\hat{\boldsymbol\Pi},
 \label{eq:dirac-even-odd-decomposition-r16}
\end{equation}
其中 $[\beta_D,\hat{\mathcal V}]=0$、$\{\beta_D,\hat{\mathcal O}\}=0$，所以 $\hat{\mathcal V}$ 保持 $P_+/Q_-$ 两个块，而 $\hat{\mathcal O}$ 只连接两个块。若研究的正能能窗满足
\begin{equation*}
 |z-m_{\mathrm e} c^2|\le\delta_\star,
 \qquad
 \|\hat{\mathcal V}\|\le V_\star,
\end{equation*}
其中 $\delta_\star$ 与 $V_\star$ 都具有能量单位 $\mathrm J$，则 $\beta_D=-1$ 块的谱与该能窗之间至少相隔
\begin{equation*}
 \Delta_-
 \equiv2m_{\mathrm e} c^2-\delta_\star-V_\star.
\end{equation*}
只要 $\Delta_->0$，谱定理对预解式给出
\begin{align}
 R_-(z)
 &\equiv\left[z-Q_-\hat H_DQ_-\right]^{-1},\\
 \|R_-(z)\|
 &=\frac{1}{\operatorname{dist}\!\left(z,\sigma(Q_-\hat H_DQ_-)\right)},\\
 &\le\frac1{\Delta_-}.
 \label{eq:dirac-negative-resolvent-bound-r16}
\end{align}
这里 $\|\cdot\|$ 是算符范数，$\sigma(\cdot)$ 表示算符谱。因此两个静止能量块之间混合的直接无量纲控制量可以取为
\begin{equation}
 \boxed{
 \epsilon_{\rm NR}
 \equiv
 \frac{\|P_+\hat{\mathcal O}Q_-\|}{\Delta_-}\ll1.}
 \label{eq:dirac-pauli-control-r16}
\end{equation}
在弱外场且典型动量为 $p_\star$ 时，$\|\hat{\mathcal O}\|\sim cp_\star$、$\Delta_-\simeq2m_{\mathrm e} c^2$，所以 $\epsilon_{\rm NR}\sim p_\star/(2m_{\mathrm e} c)$。该比值把正能投影与真实的正负能谱隙联系起来，从而给出消去负能子空间的定量条件。

受控低能分支的静态电磁哈密顿量可由 Foldy--Wouthuysen 变换逐阶分块对角化。时变外场还会从幺正变换的 $-\ii\hbar U\partial_tU^\dagger$ 项产生含 $\partial_t\mathbf E$ 的附加项，因此静态展开只适用于外场时间变化满足相应低频条件的情形。

Foldy--Wouthuysen 变换用幺正变换逐阶消去关于 $\beta_D$ 的奇算符。对静态外场展开到 $1/(m_{\mathrm e} c^2)$ 的相对论低能修正阶，分块对角的偶算符结构为
\begin{equation}
 \boxed{
 H_{\rm FW}
 =\beta_Dm_{\mathrm e} c^2+\hat{\mathcal V}
 +\frac{\beta_D\hat{\mathcal O}^2}{2m_{\mathrm e} c^2}
 -\frac{\beta_D\hat{\mathcal O}^4}{8m_{\mathrm e}^3c^6}
 -\frac{1}{8m_{\mathrm e}^2c^4}[\hat{\mathcal O},[\hat{\mathcal O},\hat{\mathcal V}]]
 +\cdots.}
 \label{eq:fw-even-operator-form-r16}
\end{equation}
这条式子先保持抽象的奇、偶算符，因此与具体电磁场分量无关；电磁结构只进入 $\hat{\mathcal O}^2$ 与双重对易子。定义无量纲四分量自旋矩阵
\begin{equation*}
 \boldsymbol\Sigma
 \equiv
 \begin{pmatrix}
 \boldsymbol\sigma_{\rm P}&0\\
 0&\boldsymbol\sigma_{\rm P}
 \end{pmatrix},
 \qquad
 \mathbf S=\frac{\hbar}{2}\boldsymbol\Sigma,
\end{equation*}
其中 $\boldsymbol\sigma_{\rm P}$ 是 Pauli 矩阵向量。利用动力学动量曲率，奇算符平方化为
\begin{equation}
 \boxed{
 \hat{\mathcal O}^2
 =c^2\left(
 \boldsymbol\Pi^2-q_{\mathrm e}\hbar\boldsymbol\Sigma\cdot\mathbf B
 \right).}
 \label{eq:fw-O2-r16}
\end{equation}
第一项给出动能，第二项已经包含自旋--磁场耦合。对静态标势 $\hat{\mathcal V}=q_{\mathrm e}\phi$，双重对易子进一步给出
\begin{equation}
 \boxed{
 \begin{aligned}
 [\hat{\mathcal O},[\hat{\mathcal O},\hat{\mathcal V}]]
 ={}&q_{\mathrm e}\hbar^2c^2\,\boldsymbol\nabla\cdot\mathbf E\\
 &+q_{\mathrm e}\hbar c^2\,
 \boldsymbol\Sigma\cdot
 \left(
 \mathbf E\times\boldsymbol\Pi
 -\boldsymbol\Pi\times\mathbf E
 \right).
 \end{aligned}}
 \label{eq:fw-double-commutator-r16}
\end{equation}
其中散度项来自 $[\Pi_i,E_i]$，叉乘项来自 Pauli 矩阵乘法中的反对称部分。把这两个电磁恒等式代回 \eqnrefs{eq:fw-even-operator-form-r16} 并投影到低能正能块，得到
\begin{equation}
 \boxed{
 \begin{aligned}
 \hat H_{+}
 ={}&m_{\mathrm e} c^2+q_{\mathrm e}\phi
 +\frac{\hat{\boldsymbol\Pi}^2}{2m_{\mathrm e}}
 -\frac{q_{\mathrm e}\hbar}{2m_{\mathrm e}}\boldsymbol\sigma_{\rm P}\cdot\mathbf B
 -\frac{\left(\hat{\boldsymbol\Pi}^2-q_{\mathrm e}\hbar\boldsymbol\sigma_{\rm P}\cdot\mathbf B\right)^2}{8m_{\mathrm e}^3c^2}\\
 &-\frac{q_{\mathrm e}\hbar^2}{8m_{\mathrm e}^2c^2}\boldsymbol\nabla\cdot\mathbf E
 -\frac{q_{\mathrm e}\hbar}{8m_{\mathrm e}^2c^2}
 \boldsymbol\sigma_{\rm P}\cdot
 \left(\mathbf E\times\hat{\boldsymbol\Pi}
 -\hat{\boldsymbol\Pi}\times\mathbf E\right)+\cdots.
 \end{aligned}}
 \label{eq:fw-pauli}
\end{equation}
依次可辨认出静质量、标势、非相对论动能、磁矩、相对论动能与弱场联合修正、Darwin 项和自旋--轨道项。这里四次奇算符项保留为完整算符平方；由于 $\hat{\boldsymbol\Pi}^2$ 与位置依赖的 $\mathbf B$ 一般不对易，不能在没有额外弱场计数时把它无条件替换成 $\hat{\boldsymbol\Pi}^4$。若进一步满足 $|q_{\mathrm e}|\hbar B_\star\ll p_\star^2$，才可在这一项内舍去磁场交叉项并恢复常见的 $-\hat{\boldsymbol\Pi}^4/(8m_{\mathrm e}^3c^2)$ 形式。所有这些项仍来自同一个 Dirac Hamiltonian 的分块对角化。
为把上述低能展开写成定量适用条件，令 $p_\star$ 是电子典型动力学动量，$\phi_\star,E_\star,B_\star$ 是外场在作用区的典型幅值，$L_\star$ 与 $\omega_\star$ 分别是其最短空间变化尺度和最高相关角频率，并定义 $\lambda_C\equiv\hbar/(m_{\mathrm e} c)$。一组足够的无量纲控制量是
\begin{equation}
 \max\!\left\{
 \frac{p_\star}{m_{\mathrm e} c},
 \frac{|q_{\mathrm e}|\,|\phi_\star|}{m_{\mathrm e} c^2},
 \frac{\hbar\omega_\star}{m_{\mathrm e} c^2},
 \frac{\lambda_C}{L_\star},
 \frac{|q_{\mathrm e}|\hbar B_\star}{m_{\mathrm e}^2c^2},
 \frac{|q_{\mathrm e}|E_\star\lambda_C}{m_{\mathrm e} c^2}
 \right\}\ll1.
 \label{eq:fw-control-set-r16}
\end{equation}
其中各项分别控制动量、标势、时间变化、空间梯度、磁场能标与电场在 Compton 长度上的势降。对数十至数百 keV 电子，第一项通常并不小，因此自由电子动力学必须保留完整相对论色散；\eqnrefs{eq:fw-pauli} 只在满足 \eqnrefs{eq:dirac-pauli-control-r16,eq:fw-control-set-r16} 时作为低能特例。

\subsection{自由场与相对论运动学}\label{sec:feqo-positive-spinor}

把电磁联络置零以后，Maxwell--Dirac 理论的电子部分退化为自由 Dirac 方程
\begin{equation}
 (\ii\hbar\slashed\partial-m_{\mathrm e} c)\psi=0.
 \label{eq:free-dirac}\end{equation}
令 $A_\mu=0$ 得到含场方程的自由极限；它同时是自由传播、散射外线态与反冲运动学的共同母方程。对正频平面波
\begin{equation*}
 \psi(x)=u_s(p)\ee^{-\ii p\cdot x/\hbar},
 \qquad p^0=E_{\mathbf p}/c>0,
\end{equation*}
\eqnrefs{eq:free-dirac} 给出
\begin{equation}
 (\slashed p-m_{\mathrm e} c)u_s(p)=0,
 \label{eq:u-dirac}
\end{equation}
而负频分支写成 $v_s(p)\ee^{+\ii p\cdot x/\hbar}$，满足
\begin{equation}
 (\slashed p+m_{\mathrm e} c)v_s(p)=0.
 \label{eq:v-dirac}
\end{equation}
两类旋量都用未来指向的 $p^0>0$ 参数化；量子化后，$u$ 与 $v$ 分支分别乘电子湮灭算符与正电子产生算符。

Clifford 代数使一阶 Dirac 算符成为相对论质量壳的因式分解，因此非零自由解必须满足
\begin{equation}
 \boxed{
 p^2=m_{\mathrm e}^2c^2,
 \qquad
 E_{\mathbf p}=\sqrt{m_{\mathrm e}^2c^4+c^2\mathbf p^2}.}
 \label{eq:rel-dispersion}\end{equation}
这条精确色散关系是传播与反冲近似的共同母式；任何近轴、低反冲或二阶色散公式都必须从它展开，而不能另行假设抛物线能谱。

为固定外线旋量归一化，取两组二分量正交基 $\xi_s^{(+)}$ 与 $\xi_s^{(-)}$，分别对应正能与负能支，并满足
\[
(\xi_s^{(+)})^\dagger\xi_{s'}^{(+)}=(\xi_s^{(-)})^\dagger\xi_{s'}^{(-)}=\delta_{ss'}.
\]
在 Dirac 表示中可选
\begin{equation}
 u_s(p)=
 \sqrt{\frac{E}{c}+m_{\mathrm e} c}
 \begin{pmatrix}
 \xi_s^{(+)}\\[2mm]
 \dfrac{\boldsymbol\sigma_{\rm P}\cdot\mathbf p}{E/c+m_{\mathrm e} c}\,\xi_s^{(+)}
 \end{pmatrix},
 \label{eq:explicit-u}
\end{equation}
负能支由同一分块方程得到，只是大、小分量的位置互换：
\begin{equation}
 v_s(p)=
 \sqrt{\frac{E}{c}+m_{\mathrm e} c}
 \begin{pmatrix}
 \dfrac{\boldsymbol\sigma_{\rm P}\cdot\mathbf p}{E/c+m_{\mathrm e} c}\,\xi_s^{(-)}\\[2mm]
 \xi_s^{(-)}
 \end{pmatrix},
 \label{eq:explicit-v}
\end{equation}
并采用
$\bar u_su_{s'}=2m_{\mathrm e} c\,\delta_{ss'}$、
$\bar v_sv_{s'}=-2m_{\mathrm e} c\,\delta_{ss'}$。
散射的自旋求和使用协变旋量空间中的代数投影。由 $\slashed p^{\,2}=m_{\mathrm e}^2c^2$ 定义
\begin{equation*}
\Lambda_+(p)\equiv\frac{\slashed p+m_{\mathrm e} c}{2m_{\mathrm e} c},
\qquad
\Lambda_-(p)\equiv\frac{m_{\mathrm e} c-\slashed p}{2m_{\mathrm e} c}.
\end{equation*}
它们满足 $\Lambda_\pm^2=\Lambda_\pm$、$\Lambda_+\Lambda_-=0$ 且迹均为 $2$；$\Lambda_+$ 的像空间正是 $(\slashed p-m_{\mathrm e} c)u=0$ 的二维正能解空间，$\Lambda_-$ 的像空间正是 $(\slashed p+m_{\mathrm e} c)v=0$ 的二维负频解空间。它们满足 $\Lambda_\pm^\dagger=\gamma^0\Lambda_\pm\gamma^0$，一般不对普通正定内积自伴，因此与下面用于单电子概率的 Hamiltonian 正交投影不同。结合上面的 Dirac 归一化，得到协变完备关系
\begin{align}
 \boxed{\sum_su_s(p)\bar u_s(p)=\slashed p+m_{\mathrm e} c,}
 &\qquad \label{eq:u-complete}\\
 \boxed{\sum_sv_s(p)\bar v_s(p)=\slashed p-m_{\mathrm e} c.}
 &\qquad \label{eq:v-complete}
\end{align}
它们由任意在壳 $p$ 上的 Dirac 投影算符直接给出；在散射振幅里，外部自旋求和因此可以化成伽马矩阵迹，同时固定 Dirac 传播子的分子结构。

协变旋量用于散射外线，而实际束流由普通概率归一化的正能波包制备。两种归一化必须在进入跃迁核之前换算，因为概率投影、自旋求和与相空间使用的内积不同。

\begin{definition}[自由哈密顿量与概率归一化旋量]
令 $\beta=\gamma^0$、$\alpha_i=\gamma^0\gamma^i$，则
\begin{align}
 \{\alpha_i,\alpha_j\}&=2\delta_{ij}I_4,\qquad
 \{\alpha_i,\beta\}=0,\qquad \beta^2=I_4,
 \label{eq:feqo-dirac-anticommutators}\\
 H_0(\vb{p})&=c\boldsymbol\alpha\cdot\vb{p}+\beta m_{\mathrm e}c^2,
 \qquad H_0^2=E_{\vb{p}}^2I_4.
 \label{eq:feqo-dirac-free}
\end{align}
这里 $\alpha_i$ 与 $\beta$ 都厄米，$E_{\vb{p}}>0$。协变式和哈密顿量式之间的词典是
\begin{equation}
 \gamma^0=\beta,\qquad \gamma^i=\beta\alpha_i.
 \label{eq:feqo-gamma-dictionary}
\end{equation}
将\eqnrefs{eq:explicit-u} 的旋量缩放为 $U_s(\vb{p})=\sqrt{c/(2E_{\vb{p}})}\,u_s(p)$，便有
\begin{equation}
 U_s(\vb{p})=N_{\vb{p}}
 \begin{pmatrix}\xi_s\\ C_{\vb{p}}\xi_s\end{pmatrix},
 \qquad
 N_{\vb{p}}=\sqrt{\frac{E_{\vb{p}}+m_{\mathrm e}c^2}{2E_{\vb{p}}}},
 \qquad
 C_{\vb{p}}=\frac{c\boldsymbol\sigma_{\rm P}\cdot\vb{p}}
 {E_{\vb{p}}+m_{\mathrm e}c^2}.
 \label{eq:feqo-positive-spinor}
\end{equation}
大写 $U_s$ 在本段表示单位概率范数旋量；小写 $u_s(p)$ 始终保留\eqnrefs{eq:explicit-u} 的协变归一化。
\end{definition}

\begin{proposition}[同一动量上的正交能量分解]
在固定 $\vb{p}$ 上，定义
\begin{equation}
 \Pi_\pm(\vb{p})=\frac12\left(I_4\pm\frac{H_0(\vb{p})}{E_{\vb{p}}}\right).
 \label{eq:feqo-hamiltonian-energy-projectors}
\end{equation}
则 $\Pi_\pm^2=\Pi_\pm=\Pi_\pm^\dagger$、$\Pi_+\Pi_-=0$、$\Pi_++\Pi_-=I_4$，且 $\Tr\Pi_\pm=2$、$H_0\Pi_\pm=\pm E_{\vb{p}}\Pi_\pm$。同一三动量上的负能本征旋量可取
\begin{equation}
 W_s(\vb{p})=N_{\vb{p}}
 \begin{pmatrix}-C_{\vb{p}}\eta_s\\\eta_s\end{pmatrix},
 \qquad H_0(\vb{p})W_s(\vb{p})=-E_{\vb{p}}W_s(\vb{p}),
 \label{eq:feqo-negative-hamiltonian-spinor}
\end{equation}
其中 $\eta_s$ 是二分量单位正交基。$\{U_1,U_2,W_1,W_2\}$ 是普通内积下的完整正交基。
\end{proposition}

\begin{proof}
Dirac 本征方程的上下两块为
\begin{equation*}
 (E-m_{\mathrm e}c^2)\phi=c\boldsymbol\sigma_{\rm P}\cdot\vb{p}\,\chi,
 \qquad
 (E+m_{\mathrm e}c^2)\chi=c\boldsymbol\sigma_{\rm P}\cdot\vb{p}\,\phi.
\end{equation*}
正能根给 $\chi=C_{\vb{p}}\phi$；负能根给 $\phi=-C_{\vb{p}}\chi$。因为 $C_{\vb{p}}$ 厄米且
\begin{equation*}
 C_{\vb{p}}^2=
 \frac{E_{\vb{p}}-m_{\mathrm e}c^2}{E_{\vb{p}}+m_{\mathrm e}c^2}I_2,
 \qquad N_{\vb{p}}^2(I_2+C_{\vb{p}}^2)=I_2,
\end{equation*}
有 $U_s^\dagger U_{s'}=W_s^\dagger W_{s'}=\delta_{ss'}$。不同能量支的正交性也直接给出
\begin{equation*}
 U_s^\dagger W_{s'}=N_{\vb{p}}^2\xi_s^\dagger(-C_{\vb{p}}+C_{\vb{p}})\eta_{s'}=0.
\end{equation*}
令 $h=H_0/E_{\vb{p}}$，则 $h=h^\dagger$、$h^2=I_4$，从而 $(I_4\pm h)^2=2(I_4\pm h)$ 和 $(I_4+h)(I_4-h)=0$ 给出所有投影恒等式。$\Tr H_0=0$ 又给 $\Tr\Pi_\pm=2$。因此
\begin{equation*}
 \Pi_+=\sum_sU_sU_s^\dagger,
 \qquad \Pi_-=\sum_sW_sW_s^\dagger.
\end{equation*}
对于\eqnrefs{eq:explicit-v}，$H_0(-\vb{p})v_s(p)=-E_{\vb{p}}v_s(p)$，故 $W_s(\vb{p})=\sqrt{c/(2E_{\vb{p}})}\,v_s(p^0,-\vb{p})$。负频反粒子旋量与同一三动量的负能哈密顿量旋量不能使用同一个动量实参。
\end{proof}

\begin{proposition}[协变自旋和与概率投影的换算]
未来指向在壳动量 $p$ 对应
\begin{align}
 \sum_sU_s(\vb{p})\bar U_s(\vb{p})
 &=\Pi_+(\vb{p})\gamma^0
 =\frac{c(\slashed p+m_{\mathrm e}c)}{2E_{\vb{p}}},
 \label{eq:feqo-probability-spin-sum}\\
 \sum_sW_s(\vb{p})\bar W_s(\vb{p})
 &=\Pi_-(\vb{p})\gamma^0
 =\frac{c(\slashed{\tilde p}-m_{\mathrm e}c)}{2E_{\vb{p}}},
 \qquad \tilde p=(E_{\vb{p}}/c,-\vb{p}).
 \label{eq:feqo-negative-probability-spin-sum}
\end{align}
普通内积的正交投影为 $\Pi_\pm$；含 Dirac 共轭的自旋和并非普通正交投影。
\end{proposition}

\begin{proof}
右乘 $\gamma^0$ 时，$\alpha_i\gamma^0=-\gamma^i$，所以
\begin{equation*}
 \Pi_+\gamma^0=\frac{E_{\vb{p}}\gamma^0-c\gamma^ip_i+m_{\mathrm e}c^2I_4}{2E_{\vb{p}}},
 \qquad
 \Pi_-\gamma^0=\frac{E_{\vb{p}}\gamma^0+c\gamma^ip_i-m_{\mathrm e}c^2I_4}{2E_{\vb{p}}}.
\end{equation*}
两式分别是 $c(\slashed p+m_{\mathrm e}c)/(2E_{\vb{p}})$ 和 $c(\slashed{\tilde p}-m_{\mathrm e}c)/(2E_{\vb{p}})$。同时
\begin{equation*}
 \bar U_sU_{s'}=\frac{m_{\mathrm e}c^2}{E_{\vb{p}}}\delta_{ss'},
 \qquad
 \bar W_sW_{s'}=-\frac{m_{\mathrm e}c^2}{E_{\vb{p}}}\delta_{ss'}.
\end{equation*}
协变归一化只需乘回 $\sqrt{2E_{\vb{p}}/c}$，得到\eqnrefs{eq:u-complete,eq:v-complete}。
\end{proof}

\begin{exbox}[$200\,\mathrm{keV}$ 电子的上下分量与投影]
取动能为 $200\,\mathrm{keV}$ 的电子，$\vb{p}=p\vb{e}_z$、$m_{\mathrm e}c^2=511\,\mathrm{keV}$。求该电子正能旋量的上下分量比、基矢的正能投影概率与流速。
\textbf{解。} 总能量是 $E=711\,\mathrm{keV}$，
\begin{equation*}
 pc=\sqrt{711^2-511^2}\,\mathrm{keV}=494.4\,\mathrm{keV},
 \qquad r=\frac{pc}{E+m_{\mathrm e}c^2}=0.4046,
 \qquad N^2=\frac{1222}{1422}=0.85935.
\end{equation*}
自旋向上时 $U_\uparrow=N(1,0,r,0)^{\mathsf T}$，因此上下两块的概率分别为 $N^2=0.85935$ 和 $N^2r^2=0.14065$；下分量已占约 $14\%$。对坐标基向量 $e_1=(1,0,0,0)^{\mathsf T}$，
\begin{equation*}
 \Pi_+e_1=N^2(1,0,r,0)^{\mathsf T},
 \qquad \|\Pi_+e_1\|^2=N^2.
\end{equation*}
所以 $e_1$ 并非纯正能电子；归一化其正能部分才得到 $U_\uparrow$。同一旋量的流速为 $c^2p/E=0.6953c$，与精确色散的群速度相同。
\end{exbox}

把正能旋量叠加成波包以后，概率流由同一个 Dirac 方程确定。
\label{sec:feqo-current-and-packet}
\begin{definition}[正能单电子波包与概率流]
以普通三维 $\delta$ 归一化动量基写入射电子，
\begin{equation}
 \Psi(\vb{r},t)=\sum_s\int\frac{\dd^3p}{(2\pi\hbar)^{3/2}}
 a_s(\vb{p})U_s(\vb{p})
 \ee^{\ii(\vb{p}\cdot\vb{r}-E_{\vb{p}}t)/\hbar},
 \qquad \sum_s\int\dd^3p\,|a_s(\vb{p})|^2=1.
 \label{eq:feqo-positive-wavepacket}
\end{equation}
概率密度、概率流和电荷四流依次为
\begin{equation}
 \rho=\Psi^\dagger\Psi,\qquad
 \vb{j}_{\rm prob}=c\Psi^\dagger\boldsymbol\alpha\Psi,\qquad
 j_{\rm ch}^\mu=q_{\mathrm e}c\bar\Psi\gamma^\mu\Psi
 =(q_{\mathrm e}c\rho,q_{\mathrm e}\vb{j}_{\rm prob}).
 \label{eq:feqo-dirac-current}
\end{equation}
$[\rho]=\mathrm{m^{-3}}$、$[\vb{j}_{\rm prob}]=\mathrm{m^{-2}s^{-1}}$；进入 Maxwell 方程的是电荷流。
\end{definition}

\begin{proof}[证明（归一化、连续性与群速度）]
对位置内积先积分平面波，得到 $(2\pi\hbar)^3\delta^{(3)}(\vb{p}-\vb{q})$。在分布支撑上 $U_s^\dagger(\vb{p})U_{s'}(\vb{p})=\delta_{ss'}$，因此
\begin{equation*}
 \int\dd^3r\,\Psi^\dagger\Psi
 =\sum_s\int\dd^3p\,|a_s(\vb{p})|^2=1.
\end{equation*}
将 Dirac 哈密顿量方程与其厄米共轭相加，质量项和实标势相消，微分动能给
\begin{equation*}
 \partial_t(\Psi^\dagger\Psi)
 =-c\big[(\partial_i\Psi^\dagger)\alpha_i\Psi
 +\Psi^\dagger\alpha_i\partial_i\Psi\big]
 =-\nabla\cdot\vb{j}_{\rm prob}.
\end{equation*}
波包衰减使无穷远通量消失，所以总概率守恒。对固定动量的正能旋量，Pauli 反对易关系直接给
\begin{align*}
 U_s^\dagger\alpha_iU_s
 &=N_{\vb{p}}^2\xi_s^\dagger(\sigma_iC_{\vb{p}}+C_{\vb{p}}\sigma_i)\xi_s\\
 &=N_{\vb{p}}^2\frac{2cp_i}{E_{\vb{p}}+m_{\mathrm e}c^2}
 =\frac{cp_i}{E_{\vb{p}}}.
\end{align*}
故 $\vb{j}_{\rm prob}/\rho=c^2\vb{p}/E_{\vb{p}}=\nabla_{\vb{p}}E$。波包中不同动量还存在干涉流，不能把该单平面波比值逐点代替完整电流。
\end{proof}


\begin{exbox}[沿 $z$ 方向运动的电子]
取 $\mathbf p=p\,\mathbf e_z$、$p\ge0$，并在 $z$ 轴自旋基中选
\begin{equation*}
 \xi_\uparrow=\begin{pmatrix}1\\0\end{pmatrix},
 \qquad
 \xi_\downarrow=\begin{pmatrix}0\\1\end{pmatrix}.
\end{equation*}
记
\begin{equation*}
 A\equiv\sqrt{\frac{E}{c}+m_{\mathrm e} c},
 \qquad
 \kappa\equiv\frac{p}{E/c+m_{\mathrm e} c},
 \qquad
 E=\sqrt{m_{\mathrm e}^2c^4+c^2p^2}.
\end{equation*}
由\eqnrefs{eq:explicit-u}，两个正能旋量就是
\begin{equation}
 u_\uparrow(p)=A
 \begin{pmatrix}1\\0\\\kappa\\0\end{pmatrix},
 \qquad
 u_\downarrow(p)=A
 \begin{pmatrix}0\\1\\0\\-\kappa\end{pmatrix}.
 \label{eq:explicit-z-spinors-dp62}
\end{equation}
这里小分量与大分量之比恰为 $\kappa$。低速时 $p\ll m_{\mathrm e} c$，有 $\kappa\simeq p/(2m_{\mathrm e} c)\ll1$；超相对论时 $E\gg m_{\mathrm e} c^2$，则 $\kappa\to1$。因此所谓 Dirac 旋量的“大、小分量”不是绝对分类，而由动量尺度控制。

以 $u_\uparrow$ 为例验证归一化。Dirac 表示中 $\gamma^0=\operatorname{diag}(1,1,-1,-1)$，所以
\begin{align*}
 \bar u_\uparrow u_\uparrow
 &=A^2(1-\kappa^2)\\
 &=\left(\frac Ec+m_{\mathrm e} c\right)
 \left[1-\frac{p^2}{(E/c+m_{\mathrm e} c)^2}\right]\\
 &=\frac{(E/c+m_{\mathrm e} c)^2-p^2}{E/c+m_{\mathrm e} c}\\
 &=2m_{\mathrm e} c,
\end{align*}
最后一步使用 $E^2/c^2-p^2=m_{\mathrm e}^2c^2$。类似地
\begin{align*}
 u_\uparrow^\dagger u_\uparrow
 &=A^2(1+\kappa^2)
 =\frac{2E}{c}.
\end{align*}
因此时间分量的 Dirac 流为
\begin{equation*}
 \bar u_\uparrow\gamma^0u_\uparrow
 =u_\uparrow^\dagger u_\uparrow
 =\frac{2E}{c}.
\end{equation*}
空间 $z$ 分量用 $\alpha^3\equiv\gamma^0\gamma^3$ 计算：
\begin{align*}
 \bar u_\uparrow\gamma^3u_\uparrow
 &=u_\uparrow^\dagger\alpha^3u_\uparrow\\
 &=2A^2\kappa\\
 &=2p.
\end{align*}
$x,y$ 分量为零，故
\begin{equation}
 \boxed{\bar u_s(p)\gamma^\mu u_s(p)=2p^\mu}
 \qquad(\mathbf p=p\mathbf e_z),
 \label{eq:explicit-dirac-current-z-dp62}
\end{equation}
且结果与选取 $\uparrow$ 或 $\downarrow$ 无关。这个算例把三个彼此相关的抽象公式连在一起：旋量归一化 $\bar uu=2m_{\mathrm e} c$、概率/粒子数密度 $u^\dagger u=2E/c$，以及协变四流与四动量同向。散射迹中的 $\bar u\gamma^\mu u$ 使用已经建立的一般协变恒等式；这里的显式分量计算提供对旋量归一化、号规和四流方向的独立核验。
\end{exbox}

自由电子的群速度和波包二阶色散由同一个 $E(\mathbf p)$ 的一阶、二阶导数决定。定义 Hessian（二阶导数）矩阵
\begin{equation}
 \boxed{
 H_{ij}(\mathbf p)
 \equiv\frac{\partial^2E}{\partial p_i\partial p_j}
 =\frac{c^2}{E}
 \left(\delta_{ij}-\frac{c^2p_ip_j}{E^2}\right),}
 \label{eq:exact_dispersion_hessian}
\end{equation}
其 SI 单位为 $\mathrm{kg^{-1}}$。若中心动量沿 $z$ 轴，先由精确色散定义中心能量和群速度，
\begin{equation*}
 \mathbf p_0=p_0\mathbf e_z,
 \qquad
 E_0=\sqrt{m_{\mathrm e}^2c^4+p_0^2c^2},
 \qquad
 v_0=\frac{c^2p_0}{E_0}.
\end{equation*}
再定义中心速度比与其 Lorentz 因子
\begin{equation*}
 \beta_0\equiv\frac{v_0}{c},
 \qquad
 \gamma_0\equiv\gamma_{\rm L}(\beta_0)
 =\frac1{\sqrt{1-\beta_0^2}}
 =\frac{E_0}{m_{\mathrm e} c^2},
 \qquad
 p_0=\gamma_0m_{\mathrm e} v_0.
\end{equation*}
则 Hessian 矩阵的两个主曲率对应两个曲率质量
\begin{equation}
 \boxed{
 m_\perp\equiv\gamma_0m_{\mathrm e},
 \qquad
 m_\parallel\equiv\gamma_0^3m_{\mathrm e},}
 \qquad
 H_{xx}=H_{yy}=\frac1{m_\perp},
 \quad
 H_{zz}=\frac1{m_\parallel}.
 \label{eq:dispersion-hessian-principal-r15}
\end{equation}
$m_\perp,m_\parallel$ 定义为色散在横向与纵向方向的曲率质量；静质量仍记为 $m_{\mathrm e}$。其中纵向关系也可写成
\begin{equation}
 \left.\frac{\dd^2 E}{\dd p^2}\right|_{p_0}
 =\frac1{\gamma_0^3m_{\mathrm e}}
 =\frac1{m_\parallel}.
 \label{eq:longitudinal-mass}
\end{equation}

任意自由纵向波包仍由精确 Fourier 积分传播
\begin{equation}
 \psi(z,t)=\frac{1}{\sqrt{2\pi\hbar}}
 \int_{-\infty}^{\infty}\dd p\,
 \widetilde\psi(p,0)
 \exp\!\left[\frac{\ii}{\hbar}\bigl(pz-E(p)t\bigr)\right].
 \label{eq:exact-wavepacket}
\end{equation}
只有当动量支持集中在 $p_0$ 附近时，才令 $\delta p=p-p_0$ 并展开
\begin{equation}
 E(p)=E_0+v_0\delta p+\frac{(\delta p)^2}{2m_\parallel}
 +\frac{E^{(3)}(p_0)}{3!}(\delta p)^3+\cdots.
 \label{eq:dispersion-taylor}
\end{equation}
若波包典型宽度为 $\sigma_p$，二阶近似可由无量纲比
\begin{equation}
 \boxed{
 \epsilon_{\rm curv}\equiv
 \frac{|E^{(3)}(p_0)|\sigma_p^3/6}
 {|E''(p_0)|\sigma_p^2/2}
 =\beta_0^2\frac{\sigma_p}{p_0}\ll1},
 \label{eq:third-order-control}
\end{equation}
直接控制；这比只写 $\sigma_p\ll p_0$ 更接近真正被忽略的三阶色散误差。

\begin{figure}[htbp]
 \centering
 \includegraphics[width=.78\linewidth]{dispersion_control.png}
 \caption{动能 $K_0=200\,\mathrm{keV}$ 时的精确相对论色散与在 $p_0$ 处的一、二阶展开。横轴为 $\delta p/p_0$，纵轴为 $(E-E_0)/(m_{\mathrm e} c^2)$；二阶曲线偏离精确曲线的位置直接显示高阶色散开始重要的区域。}
 \label{fig:dispersion-control}
\end{figure}

传播之外，\eqnrefs{eq:rel-dispersion} 还固定自由电子交换能量--动量时允许的运动学区域。为把吸收与发射放在同一个母式中，先把环境中一个\emph{正频模式}的四波矢定义为
\begin{equation}
 K^\mu=\left(\frac{\omega}{c},\mathbf k\right),
 \qquad \omega>0,
\end{equation}
并用符号 $\sigma=+1$ 表示电子吸收该模式量子、$\sigma=-1$ 表示电子发射该模式量子。电子四动量的有限交换统一写成
\begin{equation}
 \boxed{P_f^\mu=P_i^\mu+\sigma\hbar K^\mu.}
 \label{eq:four-transfer}
\end{equation}
这里 $K^\mu$ 固定描述正频模式本身，$\sigma$ 单独记录能量--动量流向，因此吸收与发射共享同一组 $\omega>0$、$\mathbf k$ 模式参数。初末电子都在同一质量壳 $P_i^2=P_f^2=m_{\mathrm e}^2c^2$，因此
\begin{equation}
 \boxed{2\sigma P_i\cdot K+\hbar K^2=0.}
 \label{eq:exact-recoil-four}
\end{equation}
这是一切单量子吸收/发射运动学的母恒等式。真空实光子满足 $K^2=0$；对孤立自由电子，未来指向的类时 $P_i$ 与非零未来指向类光 $K$ 有 $P_i\!\cdot K>0$，故无论 $\sigma=+1$ 还是 $-1$ 都不能满足上式。近场、周期结构、介质或散射体能够提供额外空间动量，或者使参与交换的模式具有 $K^2\neq0$，从而解除这一运动学禁阻。

对一维纵向几何，模式波数记为 $\kappa_z$。同一母条件等价于
\begin{equation}
 \boxed{E(p_0+\sigma\hbar \kappa_z)-E(p_0)=\sigma\hbar\omega.}
 \label{eq:exact-longitudinal-match}
\end{equation}
只有在 $|\hbar\kappa_z|$ 落入正文已经建立的局部 Taylor 控制窗以后，才把精确色散展开。为使截断误差本身也是一个定义明确的量，记二阶 Taylor 余项为
\begin{equation*}
 \mathcal R_3^{(\sigma)}(\kappa_z)
 \equiv
 \frac{(\sigma\hbar\kappa_z)^3}{2}
 \int_0^1\dd s\,(1-s)^2
 E^{(3)}\!\left(p_0+s\sigma\hbar\kappa_z\right).
\end{equation*}
于是精确关系可写成
\begin{equation}
 \boxed{
 \omega
 =v_0\kappa_z+\sigma\frac{\hbar \kappa_z^2}{2m_\parallel}
 +\frac{\mathcal R_3^{(\sigma)}(\kappa_z)}{\sigma\hbar}.}
 \label{eq:recoil-phase-match}
\end{equation}
余项满足
\begin{equation*}
 \left|
 \frac{\mathcal R_3^{(\sigma)}(\kappa_z)}{\hbar}
 \right|
 \le
 \frac{\hbar^2|\kappa_z|^3}{6}
 \sup_{|p-p_0|\le\hbar|\kappa_z|}|E^{(3)}(p)|.
\end{equation*}
二阶反冲项的符号随吸收/发射改变，而最低阶同步项 $v_0\kappa_z$ 与流向无关。定义曲率项相对线性项的控制量
\begin{equation*}
 \chi_r^{(1)}\equiv
 \left|\frac{\hbar \kappa_z}{2m_\parallel v_0}\right|,
\end{equation*}
只有在 $\chi_r^{(1)}\ll1$ 时才可进一步写成常用同步条件
\begin{equation}
 \omega\simeq v_0\kappa_z.
 \label{eq:linear-sync}
\end{equation}
若连续吸收同一正频模式量子，取 $\sigma=+1$；此时已占据能量梯上相邻跃迁的精确失谐必须逐级计算：
\begin{equation}
 \boxed{
 \Delta_\ell\equiv
 \frac{E(p_0+(\ell+1)\hbar \kappa_z)-E(p_0+\ell\hbar \kappa_z)}{\hbar}-\omega.}
 \label{eq:exact-ladder-detuning}
\end{equation}
从宽能量梯约化到少能级反冲模型时，真正需要比较的是被占据区间内的 $\Delta_\ell$ 变化，而不是只检查单次交换的一个失谐。

自由 Dirac 方程同时固定三种共同结构：在壳旋量与完备关系固定 QED 外态；精确相对论色散及其 Hessian 固定自由传播；同一质量壳条件固定反冲与相位匹配。低反冲、近轴和同步近似都由这些精确母式在显式小参数下约化。

\section{Maxwell Green 张量与散射}\label{sec:feqo-green}

Maxwell--Dirac 作用量给出的是局域微分方程，但真实实验通常指定的是电流源、材料和边界条件。与其为每一种几何重新解偏微分方程，更有用的做法是先问一个一般问题：在给定线性环境中，位于 $\mathbf r'$ 的单位源会在 $\mathbf r$ 产生怎样的场？这个“源到场”的线性响应核就是推迟 Green 张量。

把频域 Maxwell 方程写成线性算符以后，Green 张量就是该算符带因果边界条件的逆。由这个逆核可以统一建立三类对象：任意给定源的场响应；材料扰动引起的散射；以及由算符对称性和能量守恒导出的互易与光学恒等式。这样，几何、材料色散和边界条件都进入同一个响应核，而不必在主理论中提前列举具体实验名称。

这些结构首先属于经典线性响应。自由电子能量交换、辐射和量子化介质都使用同一个 Green 张量接口，只是源、材料本构和量子态不同；环境量子化属于另一个独立层级。

Maxwell--Dirac 场方程已经固定局域源项与真空中的因果传播。介质、边界和开放几何并不改变 Maxwell 方程的局域来源，但会改变“给定电流以后场如何传播”的逆问题。在线性时不变环境中，这个逆问题由频域 Maxwell 算符及其推迟 Green 张量统一描述；散射、耗散和后续宏观 QED 都从同一个逆核出发。

先考虑空间局域但允许各向异性和色散的线性介质。相对介电率与相对磁导率张量分别记为 $\boldsymbol\varepsilon_r(\vb{r},\omega)$、$\boldsymbol\mu_r(\vb{r},\omega)$，二者均无量纲；$\omega$ 是角频率，单位为 $\mathrm{s^{-1}}$。采用全文时间 Fourier 约定 $\ee^{-\ii\omega t}$。消去磁场以后，电场满足的二阶微分算符定义为
\begin{equation}
 {
 \mathsf L_{\rm EM}(\omega)
 \equiv
 \nabla\times\boldsymbol\mu_r^{-1}(\vb{r},\omega)\nabla\times
 -\frac{\omega^2}{c^2}\boldsymbol\varepsilon_r(\vb{r},\omega).}
 \label{eq:Maxwell-operator-SI-r12}\end{equation}
$\mathsf L_{\rm EM}$ 的量纲为长度的负二次方。材料本构、几何和边界条件都进入这个算符；若材料响应具有空间非局域性，$\boldsymbol\varepsilon_r$ 与 $\boldsymbol\mu_r$ 应相应提升为积分核，而同一算符结构仍然成立。

\begin{definition}[推迟 Maxwell Green 张量]
给定辐射出射条件，推迟并矢 Green 张量 $\mathbf G^R(\vb{r},\vb{r}';\omega)$ 定义为 $\mathsf L_{\rm EM}$ 的逆核：
\begin{equation}
 {
 \mathsf L_{\rm EM}(\omega)
 \mathbf G^R(\vb{r},\vb{r}';\omega)
 =\mathbf I\,\delta^{(3)}(\vb{r}-\vb{r}').}
 \label{eq:GR-Maxwell-SI-r12}\end{equation}
其中 $\mathbf I$ 是三维单位并矢，$[\mathbf G^R]=\mathrm{m^{-1}}$。上标 $R$ 表示推迟边界条件：源在给定频率产生的场必须满足向外辐射而不是从无穷远流入的条件。因而 Green 张量不仅包含材料参数，还包含开放系统的物理边界条件。

\end{definition}

\begin{proof}[证明（推迟边界的时间支撑）]
被动因果材料的本构响应在复频率上半平面解析；稳定系统的推迟逆核由上半平面边界值定义。对任意空间测试源与测试场先取矩阵元，定义
\begin{equation*}
 g_R(t)=\int\frac{\dd\omega}{2\pi}\,g_R(\omega)\ee^{-\ii\omega t}.
\end{equation*}
当 $t<0$，指数在上半平面衰减。若高频增长需要处理，先取足够次频率除法或对时间测试函数涂抹；由无上半平面极点的轮廓闭合，得到 $g_R(t)=0$。所有测试矩阵元均如此，因此 Green 分布先于源的响应为零。出射条件与 $\omega+\ii0^+$ 是同一推迟逆的空间和频率表述；有不稳定增益极点时不能继续调用这一被动稳定系统论证。
\end{proof}

\begin{proposition}[电流到电场的推迟响应]
若时域外加电流密度为 $\vb{J}(\vb{r},t)$，则按全文 Fourier 约定定义的频谱电流 $\widetilde{\vb{J}}(\vb{r},\omega)$ 的单位为 $\mathrm{A\,s\,m^{-2}}$。频域 Maxwell 方程可压缩成
\begin{equation}
{
\mathsf L_{\rm EM}(\omega)\widetilde{\vb{E}}(\vb{r},\omega)
=\ii\mu_0\omega\widetilde{\vb{J}}(\vb{r},\omega).}
\label{eq:maxwell-source-operator-dp68}
\end{equation}
两条旋度方程合并后得到与具体几何无关的线性源方程。利用\eqnrefs{eq:GR-Maxwell-SI-r12} 的推迟逆核，得到电流到电场的线性响应母式
\begin{equation}
 {
 \widetilde{\vb{E}}(\vb{r},\omega)
 =\ii\mu_0\omega
 \int\dd^3 r'\,
 \mathbf G^R(\vb{r},\vb{r}';\omega)
 \cdot\widetilde{\vb{J}}(\vb{r}',\omega).}
 \label{eq:E-from-GJ-SI-r12}\end{equation}
量纲检查给 $[\mu_0\omega G^R\widetilde J\,\dd^3 r']=\mathrm{V\,s\,m^{-1}}$，与时间 Fourier 电场 $\widetilde{\vb{E}}$ 的单位一致。这个公式是后续 EELS、辐射和环境诱导自能的经典响应起点；具体电子轨迹只会改变源 $\widetilde{\vb{J}}$，不会改变环境核 $\mathbf G^R$。

\end{proposition}

入射场与结构散射可由参考算符分解。选取参考算符 $\mathsf L_0$ 及其推迟 Green 张量 $\mathbf G_0^R=\mathsf L_0^{-1}$，并定义散射势算符
\begin{equation*}
 \mathsf V_{\rm EM}\equiv\mathsf L_0-\mathsf L_{\rm EM}.
\end{equation*}
同一源在真实环境中的场满足
\begin{equation}
{
\widetilde{\vb{E}}
=\widetilde{\vb{E}}_{\rm inc}+\mathbf G_0^R\mathsf V_{\rm EM}\widetilde{\vb{E}}.}
\label{eq:maxwell-lippmann-schwinger-dp68}
\end{equation}
右边第二项表示场先由结构散射一次，再由参考环境传播；把其中的总场再次用同一方程代回，就自动产生二次、三次直到任意次数的散射。只要相应预解式存在，这个级数可重求和为电磁 $T$ 算符
\begin{equation}
 {
 \mathsf T_{\rm EM}
 =\mathsf V_{\rm EM}
 (\mathbf I-\mathbf G_0^R\mathsf V_{\rm EM})^{-1}.}
 \label{eq:maxwell-t-operator-def-r16}
\end{equation}
于是频域总场 $\widetilde{\vb{E}}=\widetilde{\vb{E}}_{\rm inc}+\widetilde{\vb{E}}_{\rm sc}$ 的散射部分为
\begin{equation}
 {
 \widetilde{\vb{E}}_{\rm sc}
 =\mathbf G_0^R\mathsf T_{\rm EM}\widetilde{\vb{E}}_{\rm inc}.}
 \label{eq:maxwell_scattered_field_t_operator}
\end{equation}
这里 $\widetilde{\vb{E}}_{\rm inc}$ 是同一外部源在参考环境 $\mathsf L_0$ 中产生的 Fourier 电场。球、圆柱、平面界面和数值离散结构的差别，只在于如何表示并求解同一个 $\mathsf T_{\rm EM}$。

若局域介质满足互易本构条件
$\boldsymbol\varepsilon_r^T=\boldsymbol\varepsilon_r$ 与
$\boldsymbol\mu_r^T=\boldsymbol\mu_r$，并对两组场采用同一辐射边界条件，则分部积分后 Maxwell 算符满足
\begin{equation}
{
\int\dd^3 r\,\vb{u}\!\cdot\!\mathsf L_{\rm EM}\vb{v}
=\int\dd^3 r\,\vb{v}\!\cdot\!\mathsf L_{\rm EM}\vb{u}.}
\label{eq:maxwell-bilinear-reciprocity-dp68}
\end{equation}
这条式子只是说：在互易介质中，交换“源场”和“测试场”不会改变这一双线性响应。把两组测试场选成点源产生的 Green 场，就得到 Lorentz 互易关系
\begin{equation}
 {
 \mathbf G^R(\vb{r},\vb{r}';\omega)
 =(\mathbf G^R)^T(\vb{r}',\vb{r};\omega).}
 \label{eq:maxwell-green-reciprocity-r16}\end{equation}
互易性约束交换源点与观测点时的张量转置；它与“无损”是不同条件，因此有吸收的互易介质仍可满足此式。

Green 张量的虚部同时记录材料吸收和出射辐射。若两个逆算符在同一完整定义域上存在，预解式代数给
\begin{equation}
 {
 \mathbf G^R-(\mathbf G^R)^\dagger
 =(\mathbf G^R)^\dagger
 (\mathsf L_{\rm EM}^\dagger-\mathsf L_{\rm EM})
 \mathbf G^R.}
 \label{eq:resolvent-imag-identity-r12}
\end{equation}
这里的伴随包含边界定义域，不能只对微分表达式取复共轭。真空外部的出射边界与入射伴随边界不同，差值含辐射表面项；其显式功率形式见\eqnrefs{eq:feqo-open-green-optical-identity}。

若 $\boldsymbol\mu_r$ 为实厄米张量，先用趋于零的被动吸收完成整个外部空间，使积分无穷远的表面项消失，再在空间积分之后取弱吸收极限，便得到体积分形式
\begin{equation}
 {
 \frac{\omega^2}{c^2}
 \int\dd^3 s\,
 \mathbf G^R(\vb{r},\vb{s};\omega)
 \cdot\operatorname{Im}\boldsymbol\varepsilon_r(\vb{s},\omega)
 \cdot[\mathbf G^R(\vb{r}',\vb{s};\omega)]^\dagger
 =\operatorname{Im}\mathbf G^R(\vb{r},\vb{r}';\omega).}
 \label{eq:green-optical-identity-SI-r12}
\end{equation}
\eqnrefs{eq:green-optical-identity-SI-r12} 中的空间积分包括弱吸收的外部完成，并以 $\operatorname{Im}_{H}G=(G-G^\dagger)/(2\ii)$ 解释算符虚部。不能先令真空外部吸收为零再只积分有限材料区域；这种顺序会漏掉辐射谱。在明确保留真空外部的表述中，应使用\eqnrefs{eq:feqo-open-green-optical-identity} 的吸收与辐射两项。存在磁损耗时，还需加入相应的磁场吸收二次型。

\subsection{真空响应与辐射通道}
\label{sec:feqo-vacuum-dyadic}
\begin{proposition}[真空并矢与纵横分解]
对 $\omega>0$ 取 $k=(\omega+\ii0^+)/c$，$R=|\vb{r}-\vb{r}'|$。真空推迟核是
\begin{equation}
 \mathbf G_0^R(\vb{r},\vb{r}';\omega)
 =\left(\mathbf I+\frac{\nabla\nabla}{k^2}\right)
 \frac{\ee^{\ii kR}}{4\pi R}.
 \label{eq:feqo-vacuum-green}
\end{equation}
其中二阶导数按分布定义。若空间 Fourier 变量为 $\vb{q}$，定义 $P_L=\vb{q}\vb{q}^{\mathsf T}/q^2$、$P_T=\mathbf I-P_L$，则
\begin{equation}
 \mathbf G_0^R(\vb{q},\omega)
 =\frac{P_T}{q^2-k^2}-\frac{P_L}{k^2}.
 \label{eq:feqo-vacuum-transverse-longitudinal}
\end{equation}
完整 Green 张量包含纵向静电响应；只有源点外的自由传播辐射满足横向条件。
\end{proposition}

\begin{proof}
标量核 $g_0=\ee^{\ii kR}/(4\pi R)$ 满足 $(\nabla^2+k^2)g_0=-\delta^{(3)}(\vb{r}-\vb{r}')$，其 $1/(4\pi)$ 系数由包围源点的小球通量固定。令 $\mathsf L_0=\nabla\times\nabla\times-k^2\mathbf I$，逐项作用于并矢：
\begin{align*}
 \mathsf L_0(\mathbf I g_0)&=\nabla\nabla g_0-\mathbf I(\nabla^2+k^2)g_0,\\
 \mathsf L_0(k^{-2}\nabla\nabla g_0)&=-\nabla\nabla g_0.
\end{align*}
相加得到 $\mathsf L_0\mathbf G_0^R=\mathbf I\delta^{(3)}$。在空间 Fourier 表象，$\mathsf L_0=(q^2-k^2)P_T-k^2P_L$，两个正交子空间分别求逆便得到\eqnrefs{eq:feqo-vacuum-transverse-longitudinal}。取散度也显示接触项的作用：
\begin{equation*}
 \nabla\cdot\mathbf G_0^R
 =\nabla(g_0+k^{-2}\nabla^2g_0)
 =-k^{-2}\nabla\delta^{(3)}(\vb{r}-\vb{r}').
\end{equation*}
因此“完整并矢在分布意义下处处无散度”并不成立。
\end{proof}

\begin{exbox}[点电偶极与真空辐射功率]
采用单频实场约定 $\vb{E}_{\rm real}=\operatorname{Re}[\vb{\mathcal E}\ee^{-\ii\omega t}]$。位于 $\vb{r}_0$ 的电偶极幅度 $\vb{d}$（单位 $\mathrm{C\,m}$）产生
\begin{equation*}
 \vb{J}=-\ii\omega\vb{d}\,\delta^{(3)}(\vb{r}-\vb{r}_0),
 \qquad \rho=-\vb{d}\cdot\nabla\delta^{(3)}(\vb{r}-\vb{r}_0),
 \qquad -\ii\omega\rho+\nabla\cdot\vb{J}=0.
\end{equation*}
求远场辐射的角分布与总功率，并用局域谱权重核对结果。
\textbf{解。} 将这个守恒源代入\eqnrefs{eq:E-from-GJ-SI-r12}，得 $\vb{\mathcal E}=\mu_0\omega^2\mathbf G_0^R\vb{d}$。在 $kR\gg1$ 的远场，
\begin{equation*}
 \mathbf G_0^R=\frac{\ee^{\ii kR}}{4\pi R}
 (\mathbf I-\hat{\vb{R}}\hat{\vb{R}}^{\mathsf T})+O(R^{-2}).
\end{equation*}
故 $\vb{d}_\perp=\vb{d}-\hat{\vb{R}}(\hat{\vb{R}}\cdot\vb{d})$ 决定辐射方向，时间平均 Poynting 通量给
\begin{equation}
 \frac{\dd P_{\rm rad}}{\dd\Omega}
 =\frac{\omega^4}{32\pi^2\varepsilon_0c^3}|\vb{d}_\perp|^2,
 \qquad
 P_{\rm rad}=\frac{\omega^4|\vb{d}|^2}{12\pi\varepsilon_0c^3}.
 \label{eq:feqo-vacuum-dipole-radiation}
\end{equation}
最后一步用了 $\int\dd\Omega\,\hat R_i\hat R_j=4\pi\delta_{ij}/3$。实场幅度中的 $1/2$ 时间平均因子必须保留；若把正频分量的幅度直接当成 $\vb{\mathcal E}$，功率会错四倍。

同一个功率还能从局域谱权重获得。真空动量核的横向部分满足
\begin{equation*}
 \operatorname{Im}\mathbf G_0^R(\vb{r}_0,\vb{r}_0;\omega)
 =\pi\int\frac{\dd^3q}{(2\pi)^3}P_T(\vb{q})\delta(q^2-k^2)
 =\frac{k}{6\pi}\mathbf I.
\end{equation*}
这里径向积分为 $k/2$，横向投影的角积分为 $8\pi\mathbf I/3$。于是
\begin{equation*}
 \frac{\mu_0\omega^3}{2}\vb{d}^*\cdot
 \operatorname{Im}\mathbf G_0^R(\vb{r}_0,\vb{r}_0;\omega)\vb{d}
 =\frac{\omega^4|\vb{d}|^2}{12\pi\varepsilon_0c^3}.
\end{equation*}
材料完全无损时这个谱权重仍非零：它记录光子流向无穷远的辐射通道。
\end{exbox}

\begin{proposition}[源功率、吸收与辐射边界]
对频率 $\omega>0$ 的紧支撑电流幅度 $\vb{J}$，记内积为 $\langle\vb{u},\vb{v}\rangle=\int\dd^3r\,\vb{u}^*\cdot\vb{v}$，并定义算符虚部 $\operatorname{Im}_{H}G=(G-G^\dagger)/(2\ii)$。源交给电磁系统的时间平均功率是
\begin{equation}
 P_{\rm src}=-\frac12\operatorname{Re}\langle\vb{J},\vb{\mathcal E}\rangle
 =\frac{\mu_0\omega}{2}\langle\vb{J},\operatorname{Im}_{H}\mathbf G^R\vb{J}\rangle
 =P_{\rm abs}+P_{\rm rad}.
 \label{eq:feqo-maxwell-power-balance}
\end{equation}
其中 $P_{\rm abs}$ 是材料内的吸收，$P_{\rm rad}$ 是流向无穷远的功率。在被动材料及出射边界下，两者均非负；有限探测孔径只收集 $P_{\rm rad}$ 的一部分。
\end{proposition}

\begin{proof}
将 $\vb{\mathcal E}=\ii\mu_0\omega\mathbf G^R\vb{J}$ 代入第一项就得到 Green 二次型。对场占据的体积 $V$ 使用
\begin{equation*}
 \nabla\cdot(\vb{\mathcal E}\times\vb{\mathcal H}^*)
 =\vb{\mathcal H}^*\cdot\nabla\times\vb{\mathcal E}
 -\vb{\mathcal E}\cdot\nabla\times\vb{\mathcal H}^*,
\end{equation*}
代入频域 Maxwell 方程并取实部，得
\begin{align*}
 P_{\rm src}
 &=\frac12\operatorname{Re}\int_{\partial V}
 (\vb{\mathcal E}\times\vb{\mathcal H}^*)\cdot\vb{n}\,\dd A\\
 &\quad+\frac{\omega\varepsilon_0}{2}\int_V\dd^3r\,
 \vb{\mathcal E}^*\cdot\operatorname{Im}_{H}\boldsymbol\varepsilon_r\vb{\mathcal E}
 +\frac{\omega\mu_0}{2}\int_V\dd^3r\,
 \vb{\mathcal H}^*\cdot\operatorname{Im}_{H}\boldsymbol\mu_r\vb{\mathcal H}.
\end{align*}
被动性使两个体积分非负。令边界扩到所有材料外的远场，出射条件使表面积分非负。有限体积中的通量则可有任意符号，不能直接解释为吸收。
\end{proof}

\begin{definition}[开放 Maxwell 逆核的辐射谱项]
在无磁损耗且外部为真空时，定义 $\mathsf R_{\rm rad}(\omega)$ 为辐射边界通道对 Green 算符虚部的贡献。其二次型由\eqnrefs{eq:feqo-maxwell-power-balance} 唯一固定为
\begin{equation}
 \langle\vb{J},\mathsf R_{\rm rad}\vb{J}\rangle
 =\frac{2P_{\rm rad}[\vb{J}]}{\mu_0\omega}\ge0.
 \label{eq:feqo-radiative-spectral-operator}
\end{equation}
因而完整光学恒等式为
\begin{equation}
 \operatorname{Im}_{H}\mathbf G^R
 =\frac{\omega^2}{c^2}(\mathbf G^R)^\dagger
 \operatorname{Im}_{H}\boldsymbol\varepsilon_r\mathbf G^R
 +\mathsf R_{\rm rad}.
 \label{eq:feqo-open-green-optical-identity}
\end{equation}
在封闭无辐射边界下，或先让整个外部空间带上趋于零的吸收并在积分后取极限时，辐射项才可合并到体吸收表达式中。
\end{definition}

\begin{proof}
材料吸收满足
\begin{equation*}
 P_{\rm abs}=\frac{\omega\varepsilon_0}{2}\|\boldsymbol\varepsilon_{r,I}^{1/2}
 \vb{\mathcal E}\|^2
 =\frac{\mu_0\omega}{2}\left\langle\vb{J},
 \frac{\omega^2}{c^2}(\mathbf G^R)^\dagger
 \boldsymbol\varepsilon_{r,I}\mathbf G^R\vb{J}\right\rangle.
\end{equation*}
这里 $\boldsymbol\varepsilon_{r,I}=\operatorname{Im}_{H}\boldsymbol\varepsilon_r$，且用了 $\mu_0\varepsilon_0c^2=1$。由功率守恒对任意紧支撑源成立，再以极化恒等式恢复双线性矩阵元，便得到算符恒等式。真空中体吸收为零，\eqnrefs{eq:feqo-vacuum-dipole-radiation} 的非零辐射功率完全来自 $\mathsf R_{\rm rad}$。
\end{proof}





\begin{lemma}[协变导数的交换子]
对光滑规范势及光滑截面，规范协变导数的交换子为曲率在截面上的作用，带电场的具体系数按\eqnrefs{eq:covariant-derivative}固定，结果为\eqnrefs{eq:covariant-commutator}。
\end{lemma}

\begin{proof}
先把协变性要求写成对任意测试旋量都成立的算符等式
\[
D'_\mu\,U_\chi=U_\chi D_\mu,
\qquad
U_\chi(x)=\ee^{\ii q_{\mathrm e}\chi(x)/\hbar}.
\]
若局域联络仍取一阶形式
\[
D'_\mu=\partial_\mu+\frac{\ii q_{\mathrm e}}{\hbar}A'_\mu,
\]
则作用在任意 $\psi$ 上给
\begin{align*}
D'_\mu(U_\chi\psi)
={}&U_\chi\partial_\mu\psi
+\frac{\ii q_{\mathrm e}}{\hbar}U_\chi(\partial_\mu\chi)\psi
+\frac{\ii q_{\mathrm e}}{\hbar}A'_\mu U_\chi\psi,\\
U_\chi D_\mu\psi
={}&U_\chi\partial_\mu\psi
+\frac{\ii q_{\mathrm e}}{\hbar}U_\chi A_\mu\psi.
\end{align*}
由于这两式必须对任意 $\psi$ 相等，乘法算符的系数必须逐点相等：
\[
A'_\mu+\partial_\mu\chi=A_\mu,
\]
即正文\eqnrefs{eq:gauge-A-lower}。因此在给定电荷约定和一阶局域导数结构后，$A_\mu$ 的变换律不是额外选择，而是协变性唯一固定的。

再计算曲率时，不能只把 $D_\mu$ 当成普通数相乘；应让对易子作用在任意测试场上：
\begin{align*}
[D_\mu,D_\nu]\psi
={}&\frac{\ii q_{\mathrm e}}{\hbar}
\bigl[(\partial_\mu A_\nu)-(\partial_\nu A_\mu)\bigr]\psi\\
&+\left(\frac{\ii q_{\mathrm e}}{\hbar}\right)^2
(A_\mu A_\nu-A_\nu A_\mu)\psi.
\end{align*}
所有 $A\,\partial\psi$ 交叉项成对消失。对阿贝尔 $U(1)$，$A_\mu$ 是普通乘法函数，所以第二行也消失，留下
\[
[D_\mu,D_\nu]\psi
=\frac{\ii q_{\mathrm e}}{\hbar}F_{\mu\nu}\psi,
\qquad
F_{\mu\nu}=\partial_\mu A_\nu-\partial_\nu A_\mu.
\]
最后代入 $A'_\mu=A_\mu-\partial_\mu\chi$，有
\[
F'_{\mu\nu}-F_{\mu\nu}
=-\partial_\mu\partial_\nu\chi
+\partial_\nu\partial_\mu\chi=0,
\]
其中只用了光滑标量函数的混合偏导可交换。这同时解释了为什么阿贝尔场强是联络的一阶反对称导数，以及为什么非阿贝尔推广中会多出联络自身的对易子项。
\end{proof}

\begin{lemma}[动力学动量的交换子]
在本章 SI 制最小耦合约定下，$\pi_i=-\ii\hbar\partial_i-qA_i$ 满足 $[\pi_i,\pi_j]=\ii q\hbar\epsilon_{ijk}B_k$，其中 $\vb{B}=\nabla\times\vb{A}$。
\end{lemma}

\begin{proof}
正文\eqnrefs{eq:kinetic-momentum} 已写出动力学动量与正则动量、矢势之间的关系。不同空间分量的非对易性因此完全来自 $[\hat p_i,A_j]$；把位置表象的基本对易关系逐项展开，再用 $\vb{B}=\vb{\nabla}\times\vb{A}$ 收集反对称导数，就得到正文\eqnrefs{eq:pi-commutator}。


由 $\hat\Pi_i=\hat p_i-q_{\mathrm e} A_i(\vb{r},t)$ 与 $[\hat p_i,f(\vb{r},t)]=-\ii\hbar\partial_i f$，逐项展开
\begin{align*}
 [\hat\Pi_i,\hat\Pi_j]
 &=[\hat p_i-q_{\mathrm e} A_i,\hat p_j-q_{\mathrm e} A_j]\\
 &=[\hat p_i,\hat p_j]
 -q_{\mathrm e}[\hat p_i,A_j]
 -q_{\mathrm e}[A_i,\hat p_j]
 +q_{\mathrm e}^2[A_i,A_j].
\end{align*}
坐标表象中的普通阿贝尔势是乘法算符，因此 $[A_i,A_j]=0$，而正则动量满足 $[\hat p_i,\hat p_j]=0$。剩余两项为
\begin{align*}
 [\hat\Pi_i,\hat\Pi_j]
 &=-q_{\mathrm e}(-\ii\hbar\partial_iA_j)
 -q_{\mathrm e}(+\ii\hbar\partial_jA_i)\\
 &=\ii q_{\mathrm e}\hbar(\partial_iA_j-\partial_jA_i).
\end{align*}
利用 $B_k=\epsilon_{k\ell m}\partial_\ell A_m$ 以及
\begin{align*}
 \epsilon_{ijk}B_k
 &=\epsilon_{ijk}\epsilon_{k\ell m}\partial_\ell A_m\\
 &=(\delta_{i\ell}\delta_{jm}-\delta_{im}\delta_{j\ell})
 \partial_\ell A_m\\
 &=\partial_iA_j-\partial_jA_i,
\end{align*}
得到主文公式。
\end{proof}

\begin{proposition}[最小耦合 Dirac 算子的平方]
在光滑电磁背景下，利用 Clifford 关系平方 Dirac 算子，得到二阶协变波动算子及由 $\sigma^{\mu\nu}F_{\mu\nu}$ 给出的自旋项，完整系数见\eqnrefs{eq:dirac-squared}。此式为微分表达式恒等式，闭算子定义域另行指定。
\end{proposition}

\begin{proof}
正文\eqnrefs{eq:dirac-em} 是含任意时空依赖电磁场的一阶 Dirac 方程。左乘符号相反的一阶因子后，协变导数乘积的对称部分给出 $D_\mu D^\mu$，反对称部分则通过 $[D_\mu,D_\nu]$ 化成场强并与 $\sigma^{\mu\nu}$ 收缩，从而得到正文\eqnrefs{eq:dirac-squared}。


将 \eqnrefs{eq:dirac-em} 除以 $c$，
\begin{equation*}
 (\ii\hbar\slashed D-m_{\mathrm e} c)\psi=0.
\end{equation*}
左乘共轭的一阶因子 $\ii\hbar\slashed D+m_{\mathrm e} c$：
\begin{align*}
 0
 &=(\ii\hbar\slashed D+m_{\mathrm e} c)(\ii\hbar\slashed D-m_{\mathrm e} c)\psi\\
 &=\left[(\ii\hbar)^2\gamma^\mu\gamma^\nu D_\mu D_\nu-m_{\mathrm e}^2c^2\right]\psi\\
 &=\left[-\hbar^2\gamma^\mu\gamma^\nu D_\mu D_\nu-m_{\mathrm e}^2c^2\right]\psi.
\end{align*}
由 $\sigma^{\mu\nu}=\ii[\gamma^\mu,\gamma^\nu]/2$ 与 Clifford 代数，
\begin{align*}
 \gamma^\mu\gamma^\nu
 &=\frac12\{\gamma^\mu,\gamma^\nu\}
 +\frac12[\gamma^\mu,\gamma^\nu]\\
 &=\eta^{\mu\nu}\id()-\ii\sigma^{\mu\nu}.
\end{align*}
因此
\begin{align*}
 \gamma^\mu\gamma^\nu D_\mu D_\nu
 &=D_\mu D^\mu-\ii\sigma^{\mu\nu}D_\mu D_\nu.
\end{align*}
由于 $\sigma^{\mu\nu}$ 对 $\mu,\nu$ 反对称，$D_\mu D_\nu$ 的对称部分与它收缩为零：
\begin{align*}
 \sigma^{\mu\nu}D_\mu D_\nu
 &=\frac12\sigma^{\mu\nu}
 (D_\mu D_\nu-D_\nu D_\mu)\\
 &=\frac12\sigma^{\mu\nu}[D_\mu,D_\nu]\\
 &=\frac{\ii q_{\mathrm e}}{2\hbar}\sigma^{\mu\nu}F_{\mu\nu}.
\end{align*}
于是
\begin{align*}
 \gamma^\mu\gamma^\nu D_\mu D_\nu
 &=D_\mu D^\mu
 +\frac{q_{\mathrm e}}{2\hbar}\sigma^{\mu\nu}F_{\mu\nu}.
\end{align*}
代回二阶方程并整体乘以 $-1$：
\begin{align*}
 0
 &=\left[
 \hbar^2D_\mu D^\mu+m_{\mathrm e}^2c^2
 +\frac{q_{\mathrm e}\hbar}{2}\sigma^{\mu\nu}F_{\mu\nu}
 \right]\psi,
\end{align*}
即\eqnrefs{eq:dirac-squared}。
\end{proof}

\begin{proposition}[偶奇算符的二阶块对角化]
设 $H=\beta mc^2+\mathcal E+\mathcal O$，$[\beta,\mathcal E]=0$、$\{\beta,\mathcal O\}=0$。在本节静态背景及低能弱混合展开条件下，逐阶消去奇部得到\eqnrefs{eq:fw-even-operator-form-r16}；其适用阶数由所保留的嵌套交换子决定。
\end{proposition}

\begin{proof}
取第一步反厄米生成元

\begin{equation*}
 A_1\equiv\ii S_1=\frac{\beta_D\hat{\mathcal O}}{2m_{\mathrm e} c^2},
 \qquad U_1=\ee^{A_1}.
\end{equation*}
由于 $\{\beta_D,\hat{\mathcal O}\}=0$，有 $(\beta_D\hat{\mathcal O})^\dagger=-\beta_D\hat{\mathcal O}$，所以 $U_1$ 幺正。Baker--Campbell--Hausdorff 展开为
\begin{equation*}
 H^{(1)}=H_D+[A_1,H_D]
 +\frac12[A_1,[A_1,H_D]]
 +\frac1{3!}[A_1,[A_1,[A_1,H_D]]]+\cdots .
\end{equation*}
对静质量与奇算符的嵌套对易子逐阶计算，
\begin{align*}
 [A_1,\beta_Dm_{\mathrm e} c^2]&=-\hat{\mathcal O},\\
 [A_1,\hat{\mathcal O}]&=\frac{\beta_D\hat{\mathcal O}^2}{m_{\mathrm e} c^2},\\
 [A_1,[A_1,\beta_Dm_{\mathrm e} c^2]]&=-\frac{\beta_D\hat{\mathcal O}^2}{m_{\mathrm e} c^2},\\
 [A_1,[A_1,\hat{\mathcal O}]]&=-\frac{\hat{\mathcal O}^3}{m_{\mathrm e}^2c^4}.
\end{align*}
继续到四阶并按 BCH 系数收集，原来的 $+\hat{\mathcal O}$ 被一阶项抵消，偶部分分别留下
\begin{equation*}
 \frac{\beta_D\hat{\mathcal O}^2}{2m_{\mathrm e} c^2},
 \qquad
 -\frac{\beta_D\hat{\mathcal O}^4}{8m_{\mathrm e}^3c^6}.
\end{equation*}
偶势满足 $[\beta_D,\hat{\mathcal V}]=0$。第一层对易子为奇算符，第二层对易子为
\begin{equation*}
 [A_1,[A_1,\hat{\mathcal V}]]
 =-\frac{1}{4m_{\mathrm e}^2c^4}[\hat{\mathcal O},[\hat{\mathcal O},\hat{\mathcal V}]],
\end{equation*}
乘上 BCH 的 $1/2!$ 后得到
\begin{equation*}
 -\frac{1}{8m_{\mathrm e}^2c^4}[\hat{\mathcal O},[\hat{\mathcal O},\hat{\mathcal V}]].
\end{equation*}
第一次变换仍留下三阶奇算符
\begin{equation*}
 \hat{\mathcal O}_1
 =\frac{\beta_D[\hat{\mathcal O},\hat{\mathcal V}]}{2m_{\mathrm e} c^2}
 -\frac{\hat{\mathcal O}^3}{3m_{\mathrm e}^2c^4}+\cdots .
\end{equation*}
再取 $A_2=\beta_D\hat{\mathcal O}_1/(2m_{\mathrm e} c^2)$，其与静质量的对易子正好给 $-\hat{\mathcal O}_1$，因此这一奇块被下一步变换消去。按 $\|\hat{\mathcal O}\|/(m_{\mathrm e} c^2)=O(\epsilon)$、$\|\hat{\mathcal V}\|/(m_{\mathrm e} c^2)=O(\epsilon^2)$ 计数，第二步对偶部分的首个新修正高于当前保留阶数。于是得到正文\eqnrefs{eq:fw-even-operator-form-r16}。
\end{proof}

\begin{lemma}[Dirac 奇部的平方]
若 $\mathcal O=c\boldsymbol\alpha\cdot\vb{\pi}$，则 $\mathcal O^2=c^2(\vb{\pi}^2-q\hbar\boldsymbol\Sigma\cdot\vb{B})$，其中 $\boldsymbol\Sigma$ 是本章四分量自旋矩阵。
\end{lemma}

\begin{proof}
正文\eqnrefs{eq:pi-commutator} 已说明动力学动量的反对称部分由磁场控制。对最小耦合奇算符 $\hat{\mathcal O}=c\boldsymbol\alpha\cdot\vb{\Pi}$ 求平方时，Pauli 代数的对称部分给普通动能，反对称部分恰好读取这个动量对易子，因此磁矩项由最小耦合动量的非对易部分与 Dirac--Pauli 代数共同生成。

利用
\begin{equation*}
 \alpha_i\alpha_j
 =\delta_{ij}\id()+\ii\epsilon_{ijk}\Sigma_k,
\end{equation*}
有
\begin{align*}
 \hat{\mathcal O}^2
 &=c^2\alpha_i\alpha_j\Pi_i\Pi_j,\\
 &=c^2\left[
 \vb{\Pi}^2
 +\frac{\ii}{2}\epsilon_{ijk}\Sigma_k[\Pi_i,\Pi_j]
 \right].
\end{align*}
代入正文\eqnrefs{eq:pi-commutator}，
\begin{equation*}
 \frac{\ii}{2}\epsilon_{ijk}\Sigma_k[\Pi_i,\Pi_j]
 =-q_{\mathrm e}\hbar\boldsymbol\Sigma\cdot\vb{B},
\end{equation*}
从而得到正文\eqnrefs{eq:fw-O2-r16}。
\end{proof}

\begin{lemma}[静电势的双重交换子]
在静态光滑电磁势下，$[\mathcal O,[\mathcal O,q\phi]]$ 分解为势的 Laplace 算子项与自旋轨道项，精确微分表达式为\eqnrefs{eq:fw-double-commutator-r16}。反对易对称化保证算符在本章实势条件下形式自伴。
\end{lemma}

\begin{proof}
正文\eqnrefs{eq:dirac-even-odd-decomposition-r16} 中 $\hat{\mathcal V}=q_{\mathrm e}\phi$、$\hat{\mathcal O}=c\boldsymbol\alpha\cdot\vb{\Pi}$。对静态标势，第一层对易子只依赖 $[\Pi_i,\phi]$；再与 $\hat{\mathcal O}$ 对易一次后，Pauli 代数的对称部分产生 $\nabla\cdot\vb{E}$，反对称部分产生自旋与 $\vb{E}\times\vb{\Pi}$ 的耦合。

因为 $\vb{E}=-\vb{\nabla}\phi$，
\begin{align*}
 [\hat{\mathcal O},\hat{\mathcal V}]
 &=q_{\mathrm e} c\alpha_i[\Pi_i,\phi],\\
 &=\ii q_{\mathrm e}\hbar c\,\boldsymbol\alpha\cdot\vb{E}.
\end{align*}
于是
\begin{align*}
 [\hat{\mathcal O},[\hat{\mathcal O},\hat{\mathcal V}]]
 &=\ii q_{\mathrm e}\hbar c^2[\alpha_i\Pi_i,\alpha_jE_j],\\
 &=\ii q_{\mathrm e}\hbar c^2
 \left
 \{[\Pi_i,E_i]
 +\ii\epsilon_{ijk}\Sigma_k(\Pi_iE_j+E_j\Pi_i)
 \right\}.
\end{align*}
再用 $[\Pi_i,E_i]=-\ii\hbar\partial_iE_i$ 并整理 Levi--Civita 收缩，得到
\begin{equation*}
 [\hat{\mathcal O},[\hat{\mathcal O},\hat{\mathcal V}]]
 =q_{\mathrm e}\hbar^2c^2\vb{\nabla}\cdot\vb{E}
 +q_{\mathrm e}\hbar c^2\boldsymbol\Sigma\cdot
 \left(\vb{E}\times\vb{\Pi}-\vb{\Pi}\times\vb{E}\right),
\end{equation*}
即正文\eqnrefs{eq:fw-double-commutator-r16}。
\end{proof}

\begin{proposition}[Dirac 模型的 Pauli 极限]
在低动量、静态缓变电磁背景及本节 Foldy--Wouthuysen 截断条件下，正能块的有效哈密顿量为\eqnrefs{eq:fw-pauli}，同时包含动力学能、Zeeman、自旋轨道和 Darwin 项。它是带明确截断阶的模型结论。
\end{proposition}

\begin{proof}
令
\[
P_+\equiv\frac{\id()+\beta_D}{2}.
\]
因为 $P_+\beta_D P_+=P_+$，而 $\boldsymbol\Sigma$ 在正、负能静止块内分别退化成同一组 Pauli 矩阵，所以
\[
P_+\boldsymbol\Sigma P_+
=\boldsymbol\sigma_{\rm P}P_+.
\]
对任意只作用在正能块内的态，正文\eqnrefs{eq:fw-O2-r16} 因而给出
\[
P_+\hat{\mathcal O}^2P_+
=c^2\Bigl(
\hat{\vb{\Pi}}^2
-q_{\mathrm e}\hbar\boldsymbol\sigma_{\rm P}\cdot\vb{B}
\Bigr)P_+.
\]
四次奇算符不能再单独把磁场项丢掉。由 $\hat{\mathcal O}^4=(\hat{\mathcal O}^2)^2$，有严格的算符恒等式
\begin{align*}
P_+\hat{\mathcal O}^4P_+
={}&c^4
\Bigl(
\hat{\vb{\Pi}}^2
-q_{\mathrm e}\hbar\boldsymbol\sigma_{\rm P}\cdot\vb{B}
\Bigr)^2P_+\\
={}&c^4\Bigl[
\hat{\vb{\Pi}}^4
-q_{\mathrm e}\hbar\bigl\{
\hat{\vb{\Pi}}^2,
\boldsymbol\sigma_{\rm P}\cdot\vb{B}
\bigr\}
+q_{\mathrm e}^2\hbar^2
(\boldsymbol\sigma_{\rm P}\cdot\vb{B})^2
\Bigr]P_+,
\end{align*}
其中反对易子保留了 $\hat{\vb{\Pi}}$ 与非均匀磁场不对易的信息。把这两式以及正文\eqnrefs{eq:fw-double-commutator-r16} 代入\eqnrefs{eq:fw-even-operator-form-r16}，得到
\begin{align*}
P_+H_{\rm FW}P_+
={}&\Biggl[
m_{\mathrm e} c^2+q_{\mathrm e}\phi
+\frac{\hat{\vb{\Pi}}^2-q_{\mathrm e}\hbar\boldsymbol\sigma_{\rm P}\cdot\vb{B}}{2m_{\mathrm e}}\\
&-\frac{\left(\hat{\vb{\Pi}}^2-q_{\mathrm e}\hbar\boldsymbol\sigma_{\rm P}\cdot\vb{B}\right)^2}{8m_{\mathrm e}^3c^2}
-\frac{q_{\mathrm e}\hbar^2}{8m_{\mathrm e}^2c^2}\vb{\nabla}\cdot\vb{E}\\
&-\frac{q_{\mathrm e}\hbar}{8m_{\mathrm e}^2c^2}
\boldsymbol\sigma_{\rm P}\cdot
\left(
\vb{E}\times\hat{\vb{\Pi}}
-\hat{\vb{\Pi}}\times\vb{E}
\right)
\Biggr]P_+ + P_+R_{\rm FW}P_+ .
\end{align*}
$R_{\rm FW}$ 收集前一推导中没有保留的更高阶 Foldy--Wouthuysen 项。去掉末尾的投影算符，就是正文\eqnrefs{eq:fw-pauli}。若还要使用更常见的纯 $\Pi^4$ 相对论动能修正，则必须另外控制
\[
\epsilon_{B/p}
\equiv
\frac{|q_{\mathrm e}|\hbar B_\star}{p_\star^2}\ll1;
\]
这时四次项中的反对易子与 $B^2$ 项才相对于 $\Pi^4$ 受控变小。这个附加条件与正负能块解耦条件\eqnrefs{eq:dirac-pauli-control-r16} 是不同的近似。
\end{proof}

\begin{proposition}[自由 Dirac 方程的质量壳]
设 $u\ne0$ 满足 $(c\boldsymbol\alpha\cdot\vb{p}+\beta mc^2)u=Eu$。Clifford 关系给出 $E^2=c^2\vb{p}^2+m^2c^4$，两支能量为 $\pm\sqrt{c^2\vb{p}^2+m^2c^4}$。
\end{proposition}

\begin{proof}
对正频旋量左乘 $\slashed p+m_{\mathrm e} c$：

\begin{align*}
 (\slashed p+m_{\mathrm e} c)(\slashed p-m_{\mathrm e} c)u_s
 &=\left(\gamma^\mu\gamma^\nu p_\mu p_\nu-m_{\mathrm e}^2c^2\right)u_s.
\end{align*}
把矩阵乘积分成对称与反对称两部分，
\begin{align*}
 \gamma^\mu\gamma^\nu p_\mu p_\nu
 &=\frac12\{\gamma^\mu,\gamma^\nu\}p_\mu p_\nu
 +\frac12[\gamma^\mu,\gamma^\nu]p_\mu p_\nu\\
 &=\eta^{\mu\nu}p_\mu p_\nu\,\id(4)
 =p^2\id(4),
\end{align*}
因为 $p_\mu p_\nu$ 对称，而 $[\gamma^\mu,\gamma^\nu]$ 反对称。于是
\begin{align*}
 (p^2-m_{\mathrm e}^2c^2)u_s&=0.
\end{align*}
非零旋量要求
\begin{align*}
 p^2&=m_{\mathrm e}^2c^2,\\
 \frac{E_{\vb{p}}^2}{c^2}-\vb{p}^2&=m_{\mathrm e}^2c^2,\\
 E_{\vb{p}}&=\sqrt{m_{\mathrm e}^2c^4+c^2\vb{p}^2},
\end{align*}
其中最后选择正能分支。对正文\eqnrefs{eq:v-dirac} 左乘 $\slashed p-m_{\mathrm e} c$ 得到同一二阶条件，因此电子与正电子外线共享同一正能质量壳参数化。
\end{proof}

\begin{proposition}[自由负能旋量的构造]
对 $m_{\mathrm e}>0$，固定动量的自由 Dirac 哈密顿量的负能本征空间维数为二。按本章的二分量自旋基与归一化构造旋量，可得\eqnrefs{eq:explicit-v}。本章 $v_s(p)$ 的负频参数使用正能四动量 $p$，故它属于 $H(-\vb{p})$ 的负能空间；普通厄米内积的正交关系为 $u_s(p)^\dagger v_{s'}(E_{\vb{p}}/c,-\vb{p})=0$。
\end{proposition}

\begin{proof}
在 Dirac 表示中把四分量旋量写成两个二分量块。对正频支令
\[
u_s(p)=\begin{pmatrix}\phi_s\\ \chi_s\end{pmatrix}.
\]
正文\eqnrefs{eq:u-dirac} 给出
\begin{align*}
(E-m_{\mathrm e} c^2)\phi_s
-c\boldsymbol\sigma_{\rm P}\!\cdot\!\vb{p}\,\chi_s&=0,\\
c\boldsymbol\sigma_{\rm P}\!\cdot\!\vb{p}\,\phi_s
-(E+m_{\mathrm e} c^2)\chi_s&=0.
\end{align*}
第二式先确定下分量：
\[
\chi_s
=\frac{c\boldsymbol\sigma_{\rm P}\cdot\vb{p}}{E+m_{\mathrm e} c^2}\phi_s
=\frac{\boldsymbol\sigma_{\rm P}\cdot\vb{p}}{E/c+m_{\mathrm e} c}\phi_s.
\]
取
$\phi_s=\sqrt{E/c+m_{\mathrm e} c}\,\xi_s^{(+)}$，便得到正文\eqnrefs{eq:explicit-u}。用质量壳 $E^2=m_{\mathrm e}^2c^4+c^2\vb{p}^2$ 检查归一化：
\begin{align*}
\bar u_su_{s'}
&=\left(\frac Ec+m_{\mathrm e} c\right)
\left[
1-\frac{\vb{p}^2}{(E/c+m_{\mathrm e} c)^2}
\right]\delta_{ss'}\\
&=2m_{\mathrm e} c\,\delta_{ss'}.
\end{align*}

对负频支令
\[
v_s(p)=\begin{pmatrix}\phi'_s\\ \chi'_s\end{pmatrix}.
\]
正文\eqnrefs{eq:v-dirac} 的两个分量方程为
\begin{align*}
(E+m_{\mathrm e} c^2)\phi'_s
-c\boldsymbol\sigma_{\rm P}\!\cdot\!\vb{p}\,\chi'_s&=0,\\
c\boldsymbol\sigma_{\rm P}\!\cdot\!\vb{p}\,\phi'_s
-(E-m_{\mathrm e} c^2)\chi'_s&=0.
\end{align*}
第一式给
\[
\phi'_s
=\frac{\boldsymbol\sigma_{\rm P}\cdot\vb{p}}{E/c+m_{\mathrm e} c}\chi'_s.
\]
取
$\chi'_s=\sqrt{E/c+m_{\mathrm e} c}\,\xi_s^{(-)}$，即得到正文\eqnrefs{eq:explicit-v}。同样利用质量壳，
\begin{align*}
\bar v_sv_{s'}
&=\left(\frac Ec+m_{\mathrm e} c\right)
\left[
\frac{\vb{p}^2}{(E/c+m_{\mathrm e} c)^2}-1
\right]\delta_{ss'}\\
&=-2m_{\mathrm e} c\,\delta_{ss'}.
\end{align*}
最后将 $v$ 的空间动量反向，记 $K=\boldsymbol\sigma_{\rm P}\cdot\vb{p}/(E/c+m_{\mathrm e}c)$，则 $K^\dagger=K$。两旋量的分块正比于 $(\xi,K\xi)$ 与 $(-K\eta,\eta)$，因此
\[
u_s(p)^\dagger v_{s'}(E/c,-\vb{p})
=\left(E/c+m_{\mathrm e}c\right)(-\xi_s^\dagger K\eta_{s'}+\xi_s^\dagger K\eta_{s'})=0.
\]
\end{proof}

\begin{proposition}[自由旋量的完备关系]
设 $p^2=m_{\mathrm e}^2c^2$、$p^0>0$，采用本章 $u,v$ 的归一化，则
\begin{align*}
\Lambda_+(p)&=\frac{\sum_su_s(p)\bar u_s(p)}{2m_{\mathrm e}c}
=\frac{\slashed p+m_{\mathrm e}c}{2m_{\mathrm e}c},\\
\Lambda_-(p)&=-\frac{\sum_sv_s(p)\bar v_s(p)}{2m_{\mathrm e}c}
=\frac{m_{\mathrm e}c-\slashed p}{2m_{\mathrm e}c}.
\end{align*}
它们是 $\slashed p$ 的两个协变代数投影，满足 $\Lambda_\pm^2=\Lambda_\pm$、$\Lambda_+\Lambda_-=0$、$\Lambda_++\Lambda_-=I_4$。它们在普通厄米内积下通常是斜投影；同一动量的哈密顿量的能量正交投影须另行定义。未归一化的自旋求和公式即\eqnrefs{eq:u-complete}与\eqnrefs{eq:v-complete}。
\end{proposition}

\begin{proof}
质量壳给 $\slashed p^{\,2}=m_{\mathrm e}^2c^2$，因此
\begin{align*}
\Lambda_+^2
&=\frac{(\slashed p+m_{\mathrm e} c)^2}{4m_{\mathrm e}^2c^2}
=\frac{\slashed p+m_{\mathrm e} c}{2m_{\mathrm e} c}=\Lambda_+,\\
\Lambda_-^2
&=\frac{(m_{\mathrm e} c-\slashed p)^2}{4m_{\mathrm e}^2c^2}
=\Lambda_-.
\end{align*}
并且
\begin{equation*}
(\slashed p-m_{\mathrm e} c)\Lambda_+=0,
\qquad
(\slashed p+m_{\mathrm e} c)\Lambda_-=0.
\end{equation*}
由于 $\operatorname{tr}\slashed p=0$，两投影的迹均为 $2$，所以其像空间分别与正能、负频 Dirac 解空间完全相同。再由伴随 Dirac 方程 $\bar u_s\slashed p=m_{\mathrm e}c\bar u_s$ 及 $\slashed p v_{s'}=-m_{\mathrm e}cv_{s'}$ 得 $\bar u_sv_{s'}=0$；同理 $\bar v_su_{s'}=0$。因此归一化基求和在一个特征空间上恒等，在另一个特征空间上为零，与相应 $\Lambda_\pm$ 的作用完全相同。由 $\bar u_su_{s'}=2m_{\mathrm e} c\,\delta_{ss'}$，正能解空间上的恒等投影为
\begin{equation*}
\Lambda_+=\frac1{2m_{\mathrm e} c}\sum_su_s\bar u_s,
\end{equation*}
故得到\eqnrefs{eq:u-complete}。负频基满足 $\bar v_sv_{s'}=-2m_{\mathrm e} c\,\delta_{ss'}$，于是
\begin{equation*}
\Lambda_-=-\frac1{2m_{\mathrm e} c}\sum_sv_s\bar v_s,
\end{equation*}
等价于\eqnrefs{eq:v-complete}。
\end{proof}

\begin{proposition}[电子色散 Taylor 截断的控制]
在固定中心动量的有限邻域内，精确自由色散 $E(p)=\sqrt{m^2c^4+c^2p^2}$ 的二阶近似误差由三阶导数的上确界乘 $|\Delta p|^3/6$ 控制；传播时还须使该能量误差乘作用时间除以 $\hbar$ 足够小。该判据为\eqnrefs{eq:third-order-control}。
\end{proposition}

\begin{proof}
为得到以动量宽度 $\sigma_p$ 表示的局部截断判据，取工作窗边界 $|\delta p|=\sigma_p$；正文\eqnrefs{eq:dispersion-taylor} 的最后保留项和第一个舍弃项分别为
\[
\Delta E_2=\frac12E''(p_0)\sigma_p^2,
\qquad
\Delta E_3=\frac16E^{(3)}(p_0)\sigma_p^3.
\]
先从精确色散 $E(p)=\sqrt{m_{\mathrm e}^2c^4+p^2c^2}$ 计算所需导数：
\begin{align*}
E'(p)&=\frac{pc^2}{E(p)},\\
E''(p)&=\frac{m_{\mathrm e}^2c^6}{E(p)^3},\\
E^{(3)}(p)
&=-\frac{3m_{\mathrm e}^2c^8p}{E(p)^5}.
\end{align*}
在 $p=p_0$ 处用 $E_0=\gamma_0m_{\mathrm e} c^2$ 与 $p_0c/E_0=\beta_0$，两项之比为
\begin{align*}
\epsilon_{\rm curv}
&=\left|\frac{\Delta E_3}{\Delta E_2}\right|\\
&=\frac{|E^{(3)}(p_0)|\sigma_p}{3|E''(p_0)|}\\
&=\frac{p_0c^2}{E_0^2}\sigma_p\\
&=\frac{p_0^2c^2}{E_0^2}\frac{\sigma_p}{p_0}\\
&=\beta_0^2\frac{\sigma_p}{p_0}.
\end{align*}
这就是正文\eqnrefs{eq:third-order-control}。
\end{proof}

\begin{proposition}[质量壳上的反冲相位匹配]
若入射和出射电子满足同一质量壳，且按\eqnrefs{eq:four-transfer}定义传递四动量，则两条质量壳相减给出\eqnrefs{eq:recoil-phase-match}。该恒等式未预先忽略二次反冲项。
\end{proposition}

\begin{proof}
正文\eqnrefs{eq:four-transfer} 把正频模式与能量流向分开写成
\[
P_f=P_i+\sigma\hbar K,
\qquad \sigma^2=1.
\]
初态和末态都是同一静质量的真实电子，所以
\[
P_i^2=P_f^2=m_{\mathrm e}^2c^2.
\]
两式相减时静质量项完全消去，而 $\sigma^2=1$ 保留二阶反冲项：
\begin{align*}
0
&=P_f^2-P_i^2\\
&=(P_i+\sigma\hbar K)^2-P_i^2\\
&=2\sigma\hbar P_i\!\cdot K+\hbar^2K^2.
\end{align*}
除以 $\hbar$ 得到正文\eqnrefs{eq:exact-recoil-four}。因此吸收与发射的差别只出现在与 $P_i\!\cdot K$ 相乘的一阶流向符号中；模式本身仍始终取 $\omega>0$。

在一维纵向通道中，$P_f=P_i+\sigma\hbar K$ 给出
\begin{align*}
p_f&=p_0+\sigma\hbar\kappa_z,\\
E_f-E_i&=\sigma\hbar\omega,
\end{align*}
故正文\eqnrefs{eq:exact-longitudinal-match} 是四维母恒等式的能量表示。若 $|\hbar\kappa_z|$ 落在正文已经给出的 Taylor 控制窗内，对
$\delta p=\sigma\hbar\kappa_z$ 使用带积分余项的 Taylor 定理：
\begin{align*}
E(p_0+\delta p)-E(p_0)
&=E'(p_0)\delta p
+\frac{E''(p_0)}{2}\delta p^2
+\frac{\delta p^3}{2}
\int_0^1\dd s\,(1-s)^2E^{(3)}(p_0+s\delta p)\\
&=\sigma v_0\hbar\kappa_z
+\frac{\hbar^2\kappa_z^2}{2m_\parallel}
+\mathcal R_3^{(\sigma)}(\kappa_z).
\end{align*}
与 $\sigma\hbar\omega$ 比较，并利用 $1/\sigma=\sigma$，得到正文\eqnrefs{eq:recoil-phase-match}。积分余项还立刻给出
\begin{align*}
\left|\mathcal R_3^{(\sigma)}(\kappa_z)\right|
&\le
\frac{|\hbar\kappa_z|^3}{2}
\sup_{|p-p_0|\le\hbar|\kappa_z|}|E^{(3)}(p)|
\int_0^1(1-s)^2\dd s\\
&=\frac{\hbar^3|\kappa_z|^3}{6}
\sup_{|p-p_0|\le\hbar|\kappa_z|}|E^{(3)}(p)|.
\end{align*}
这里可以直接看出：线性同步项不区分吸收和发射，第一项真正区分两种流向的反冲修正出现在偶次色散项；舍弃高阶项时也有明确的误差上界，而不是只靠一个形式上的省略号。
\end{proof}

\begin{proposition}[Maxwell 散射的 T 算符]
若背景 Green 算符 $G_0$ 已固定，扰动算符为 $V$ 且 $I-G_0V$ 可逆，则 Lippmann--Schwinger 方程有唯一解，散射部分可写为 $G_0T$ 作用于入射场，其中 $T=V(I-G_0V)^{-1}$。表达式不要求 Born 展开收敛。
\end{proposition}

\begin{proof}
把背景 Maxwell Green 算符记为 $\mathbf G_0^R$，介质相对于背景的散射势记为 $\mathsf V_{\rm EM}$。正文的 Lippmann--Schwinger 方程是
\[
\widetilde{\vb{E}}
=\widetilde{\vb{E}}_{\rm inc}
+\mathbf G_0^R\mathsf V_{\rm EM}\widetilde{\vb{E}}.
\]
先不要展开级数，而把未知总场的项移到左边：
\[
(\mathbf I-\mathbf G_0^R\mathsf V_{\rm EM})\widetilde{\vb{E}}
=\widetilde{\vb{E}}_{\rm inc}.
\]
只要所考察复频率不位于该算符的谱上，就存在逆算符，故
\[
\widetilde{\vb{E}}
=(\mathbf I-\mathbf G_0^R\mathsf V_{\rm EM})^{-1}
\widetilde{\vb{E}}_{\rm inc}.
\]
散射场因此为
\begin{align*}
\widetilde{\vb{E}}_{\rm sc}
&\equiv\widetilde{\vb{E}}-\widetilde{\vb{E}}_{\rm inc}\\
&=(\mathbf I-\mathbf G_0^R\mathsf V_{\rm EM})^{-1}
\mathbf G_0^R\mathsf V_{\rm EM}\widetilde{\vb{E}}_{\rm inc}.
\end{align*}
现在用可直接验证的算符恒等式
\[
(\mathbf I-AB)^{-1}A=A(\mathbf I-BA)^{-1}.
\]
令 $A=\mathbf G_0^R$、$B=\mathsf V_{\rm EM}$，便有
\[
\widetilde{\vb{E}}_{\rm sc}
=\mathbf G_0^R
\mathsf V_{\rm EM}
(\mathbf I-\mathbf G_0^R\mathsf V_{\rm EM})^{-1}
\widetilde{\vb{E}}_{\rm inc}.
\]
因此定义
\[
{
\mathsf T_{\rm EM}
\equiv\mathsf V_{\rm EM}
(\mathbf I-\mathbf G_0^R\mathsf V_{\rm EM})^{-1}}
\]
后立即得到
\[
{
\widetilde{\vb{E}}_{\rm sc}
=\mathbf G_0^R\mathsf T_{\rm EM}\widetilde{\vb{E}}_{\rm inc}},
\]
即正文\eqnrefs{eq:maxwell_scattered_field_t_operator}。

多次散射级数只是这个精确逆算符在收敛区域内的 Neumann 展开。若某个算符范数满足 $\lVert\mathbf G_0^R\mathsf V_{\rm EM}\rVert<1$，则
\begin{align*}
(\mathbf I-\mathbf G_0^R\mathsf V_{\rm EM})^{-1}
&=\sum_{n=0}^{\infty}(\mathbf G_0^R\mathsf V_{\rm EM})^n,\\
\mathsf T_{\rm EM}
&=\mathsf V_{\rm EM}
+\mathsf V_{\rm EM}\mathbf G_0^R\mathsf V_{\rm EM}
+\mathsf V_{\rm EM}\mathbf G_0^R\mathsf V_{\rm EM}\mathbf G_0^R\mathsf V_{\rm EM}+\cdots.
\end{align*}
第 $n$ 个额外因子 $\mathbf G_0^R\mathsf V_{\rm EM}$ 对应“再传播一次、再散射一次”。即使 Neumann 级数本身不收敛，$\mathsf T_{\rm EM}$ 仍由前面的算符逆定义；其极点对应齐次方程
\[
(\mathbf I-\mathbf G_0^R\mathsf V_{\rm EM})\vb{E}=0
\]
出现非零解，也就是
$\ker(\mathbf I-\mathbf G_0^R\mathsf V_{\rm EM})\neq\{0\}$；在有限维离散化后，这才等价于相应行列式为零。它对应无外源条件下仍可存在的本征模或共振。这样 $T$ 算符同时统一了弱散射 Born 级数与强多重散射的谱结构。
\end{proof}

\begin{theorem}[互易介质的 Green 对称性]
对满足本章双线性互易条件的线性介质，在相同边界条件及辐射条件下，电磁 Green 张量满足\eqnrefs{eq:maxwell-green-reciprocity-r16}的源观察点交换与转置关系。一般磁光非互易介质不满足此前提。
\end{theorem}

\begin{proof}
正文\eqnrefs{eq:maxwell-bilinear-reciprocity-dp68} 已经把介质互易性写成 Maxwell 算符的双线性对称关系。要把它变成 Green 张量的源点交换关系，只需让两组测试场分别由两个点源产生。

取任意常矢量 \(\vb{a},\vb{b}\)，令
\begin{equation*}
\vb{u}(\vb{s})=\mathbf G^R(\vb{s},\vb{r};\omega)\cdot\vb{a},
\qquad
\vb{v}(\vb{s})=\mathbf G^R(\vb{s},\vb{r}';\omega)\cdot\vb{b}.
\end{equation*}
由 Green 方程，
\begin{equation*}
\mathsf L_{\rm EM}\vb{u}
=\vb{a}\,\delta^{(3)}(\vb{s}-\vb{r}),
\qquad
\mathsf L_{\rm EM}\vb{v}
=\vb{b}\,\delta^{(3)}(\vb{s}-\vb{r}').
\end{equation*}
代入正文\eqnrefs{eq:maxwell-bilinear-reciprocity-dp68} 后，两个空间积分分别被 \(\delta\) 函数取样：
\begin{equation*}
\vb{a}\cdot\mathbf G^R(\vb{r},\vb{r}';\omega)\cdot\vb{b}
=\vb{b}\cdot\mathbf G^R(\vb{r}',\vb{r};\omega)\cdot\vb{a}.
\end{equation*}
因为 \(\vb{a},\vb{b}\) 任意，所以
\begin{equation*}
\mathbf G^R(\vb{r},\vb{r}';\omega)
=[\mathbf G^R(\vb{r}',\vb{r};\omega)]^T,
\end{equation*}
即正文\eqnrefs{eq:maxwell-green-reciprocity-r16}。
\end{proof}

\begin{proposition}[电磁 Green 算符的光学恒等式]
对本章因果 Maxwell 算符，在共同定义域及一致辐射边界条件下，$G-G^\dagger=G(L^\dagger-L)G^\dagger$。其含耗散介质和边界通道的电磁表达为\eqnrefs{eq:green-optical-identity-SI-r12}。
\end{proposition}

\begin{proof}
先在整个外部空间加入正的被动吸收。对固定实频 $\omega>0$，场在无穷远衰减，因此旋度算符的分部积分表面项消失，$\mathsf L_{\rm EM}$ 与其伴随可按同一完整空间处理。以下代数在该正则化下完成；弱吸收极限必须放在空间积分之后。定义
\[
\mathbf G^R=\mathsf L_{\rm EM}^{-1},
\qquad
\mathbf G^{R\dagger}=(\mathsf L_{\rm EM}^{\dagger})^{-1}.
\]
对任意两个可逆算符 $A,B$，有两种等价的预解式恒等式
\[
A^{-1}-B^{-1}
=A^{-1}(B-A)B^{-1}
=B^{-1}(B-A)A^{-1}.
\]
第一式例如由
$A^{-1}BB^{-1}-A^{-1}AB^{-1}=A^{-1}-B^{-1}$ 逐项得到。取
$A=\mathsf L_{\rm EM}$、$B=\mathsf L_{\rm EM}^{\dagger}$，先用第一种次序得到
\begin{align*}
\mathbf G^R-\mathbf G^{R\dagger}
&=\mathbf G^R
(\mathsf L_{\rm EM}^{\dagger}-\mathsf L_{\rm EM})
\mathbf G^{R\dagger}.
\end{align*}
在上述外部完成下，材料和外部吸收信息集中在算符差 $\mathsf L_{\rm EM}^{\dagger}-\mathsf L_{\rm EM}$ 中。对局域电介质且磁响应无损的情形，
\[
\mathsf L_{\rm EM}
=\nabla\times\nabla\times
-\frac{\omega^2}{c^2}\boldsymbol\varepsilon_r,
\]
故
\[
\mathsf L_{\rm EM}^{\dagger}-\mathsf L_{\rm EM}
=\frac{\omega^2}{c^2}
(\boldsymbol\varepsilon_r-\boldsymbol\varepsilon_r^{\dagger})
=2\ii\frac{\omega^2}{c^2}
\operatorname{Im}\boldsymbol\varepsilon_r.
\]
代回并除以 $2\ii$：
\[
\operatorname{Im}\mathbf G^R
=\frac{\omega^2}{c^2}
\mathbf G^R
(\operatorname{Im}\boldsymbol\varepsilon_r)
\mathbf G^{R\dagger}.
\]
把算符乘积写成坐标核。若介电张量局域，
$\operatorname{Im}\boldsymbol\varepsilon_r$ 的核为
$\operatorname{Im}\boldsymbol\varepsilon_r(\vb{s},\omega)
\delta^{(3)}(\vb{s}-\vb{s}')$，于是
\[
{
\operatorname{Im}\mathbf G^R(\vb{r},\vb{r}';\omega)
=\frac{\omega^2}{c^2}\int\dd^3s\,
\mathbf G^R(\vb{r},\vb{s};\omega)
\cdot\operatorname{Im}\boldsymbol\varepsilon_r(\vb{s},\omega)
\cdot[\mathbf G^R(\vb{r}',\vb{s};\omega)]^\dagger}.
\]
这就是正文\eqnrefs{eq:green-optical-identity-SI-r12}。

该恒等式还解释了为什么 $\operatorname{Im}\mathbf G^R$ 能作为损失谱权重。为把正定性写成最直观的场能量形式，改用上面第二种等价次序，亦有
\[
\operatorname{Im}\mathbf G^R
=\frac{\omega^2}{c^2}
\mathbf G^{R\dagger}
(\operatorname{Im}\boldsymbol\varepsilon_r)
\mathbf G^R.
\]
对任意测试电流 $\vb{J}$，定义由它激发的场
\[
\vb{E}_J(\vb{s})
\equiv\int\dd^3r'\,\mathbf G^R(\vb{s},\vb{r}';\omega)\cdot\vb{J}(\vb{r}').
\]
于是左右与 $\vb{J}^*,\vb{J}$ 收缩得到
\[
\iint\dd^3r\dd^3r'\,
\vb{J}^*(\vb{r})\cdot\operatorname{Im}\mathbf G^R(\vb{r},\vb{r}')\cdot\vb{J}(\vb{r}')
=\frac{\omega^2}{c^2}\int\dd^3s\,
\vb{E}_J^*(\vb{s})\cdot
\operatorname{Im}\boldsymbol\varepsilon_r(\vb{s})\cdot\vb{E}_J(\vb{s}).
\]
对被动介质，$\operatorname{Im}\boldsymbol\varepsilon_r\succeq0$，所以右边非负。因而 Green 张量的反厄米部分不是一个形式上的“取虚部”规则，而是由材料局域耗散经过两次传播映射回源空间后得到的正半定谱核；这正是后文 EELS、LDOS 与涨落--耗散关系共同使用的数学结构。


若外部保持严格真空，入射与出射定义域不相同，应另计边界通量；该通量正是正文\eqnrefs{eq:feqo-radiative-spectral-operator} 的 $\mathsf R_{\rm rad}$。真空点偶极的 $\operatorname{Im}G_0=kI/(6\pi)$ 说明该项即使没有任何体吸收也不会消失。
\end{proof}

